% =======================================================================
%  Branch-Local Hilbert-Polya:
%  A Boundary Theorem from Riemann, Selberg, and Prime-Hub
%
%  Author: Lukas Geiger, Independent Researcher, Bernau
%  HP3 provenance/equivalence guardrail candidate v2.4 - 2026-07-11
%  Status: local post-v2.3 candidate; no upload
% =======================================================================
\documentclass[11pt]{article}

\usepackage[letterpaper,top=2.4cm,bottom=2.6cm,left=2.2cm,right=2.2cm]{geometry}
\usepackage[utf8]{inputenc}
\usepackage{cmap}
\usepackage[T1]{fontenc}
\usepackage{lmodern}
\usepackage{microtype}
\setlength{\emergencystretch}{2em}
\usepackage[english]{babel}
\usepackage{amsmath,amssymb,amsfonts,mathtools,bm}
\usepackage{graphicx,float,booktabs,array,multirow}
\usepackage{xcolor}
\usepackage[unicode=true,colorlinks=true,linkcolor=blue!60!black,citecolor=blue!60!black,urlcolor=blue!60!black]{hyperref}
\hypersetup{
  pdftitle={The Zeta Zoo: The Mathematical Side of Functional Stability Theory},
  pdfauthor={Lukas Geiger},
  pdfsubject={Branch-local Hilbert-Polya classification and source-audit revision},
  pdfkeywords={Functional Stability Theory, Zeta Zoo, Hilbert-Polya, Weil form, SGE, UCU}
}

\usepackage{amsthm}
\newtheorem{theorem}{Theorem}[section]
\newtheorem{proposition}[theorem]{Proposition}
\newtheorem{corollary}[theorem]{Corollary}
\newtheorem{lemma}[theorem]{Lemma}
\newtheorem{conjecture}[theorem]{Conjecture}
\theoremstyle{definition}
\newtheorem{definition}[theorem]{Definition}
\newtheorem{example}[theorem]{Example}
\theoremstyle{remark}
\newtheorem{remark}[theorem]{Remark}

\newcommand{\HBL}{\mathrm{HP\text{-}BL}}
\newcommand{\PSL}{\mathrm{PSL}}
\newcommand{\Spec}{\operatorname{Spec}}
\renewcommand{\Re}{\operatorname{Re}}
\newcommand{\halfline}{\tfrac{1}{2}}
\newcommand{\zer}[1]{\zeta'/\zeta(#1)}
\DeclareMathOperator{\Weil}{Weil}

\title{The Zeta Zoo\\
  \large --- The Mathematical Side of\\
  Functional Stability Theory\\[0.3em]
  \normalsize Branch-Local Hilbert-P\'olya via the Trinity
  of UCU, SGE, and Weil-Form Transversality}
\author{Lukas Geiger\\
  \small Independent Researcher, Bernau\\
  \small \href{https://orcid.org/0009-0005-7296-1534}{ORCID: 0009-0005-7296-1534}}
\date{Version 2.4 candidate -- July 2026}

\begin{document}
\maketitle

\begin{abstract}
\noindent
Functional Stability Theory (FST) is a programme that identifies a
single structural pattern---functional positivity under a gauge
constraint---as the common substrate of open problems in number
theory, mathematical physics, and cosmology \cite{RFEPFramework}.
\textbf{This paper is the mathematical side of FST.} We formalise
the trinity of meta-principles that governs the zeta-type branch of
the theory: the \emph{Universal Convexity Uniqueness} (UCU) lemma,
which turns strict convexity of the governing functional into the
unique saturation state; the \emph{Semigroup-Group Equivalence}
(SGE) conjecture, which is used here as a conjectural audit heuristic
for when the Hilbert-P\'olya property may be an algebraic invariant of
the transfer provenance; and the
branch-transversal \emph{Weil quadratic form}, which supplies gap
positivity across both sides of the SGE dichotomy. Within this
trinity, the classical picture reorganises cleanly: the
non-existence results NE-A and NE-B of the Riemann v2.1 programme
\cite{Geiger2026RHII} exclude the algebraic routes for $\zeta(s)$;
Selberg's Laplace-Beltrami operator realises the Hilbert-P\'olya
structure explicitly for $Z_\Gamma(s)$; the Spectral Zookeeper
\cite{Zookeeper2026} closes the Riemann $C^{(2)}$/MS2 step in the
CCM/Fourier branch; the conjectural Prime-Hub graph remains open as
an explicit OP5 construction under four documented structural
obstructions, including the determinant-orientation obstruction
between $1/\zeta(s)$ and eigenvalue-one events.
We classify the \emph{zeta zoo} along the proposed SGE-dichotomy and add finite
calibration checks on five test families (Riemann, Selberg, Prime-Hub,
Dedekind $\mathbb{Q}(\sqrt{-5})$, Ihara-Petersen). These checks support
the taxonomy but do not prove SGE or any global Hilbert--P\'olya
transfer. We also derive
consequences for Dirichlet $L$-functions
\cite{PaleyWienerDirichlet, DirichletAtlas}. We close with a
cryptographic reading of the trinity as the algebraic architecture
of the Bitcoin blockchain protocol, and with the narrative circle
from the Cooperative Renormalization Model (CRM) of dark energy
through FST to the abstract trinity and back via hyperbolic
geometry. The physical instantiation of the FST programme, via the
Renormalized Free Energy Principle (RFEP), is developed in the
companion paper \cite{RFEPFramework}.
\end{abstract}

\clearpage
\tableofcontents
\clearpage

% ============================================================
\section{Introduction}
\label{sec:intro}

\subsection{The Hilbert-P\'olya conjecture and its implicit branch
choice}
\label{sec:hp-implicit-branch}

The Hilbert-P\'olya conjecture asks whether the non-trivial zeros of
the Riemann zeta function correspond to the eigenvalues of a
self-adjoint operator. Its earliest documented origin is a 1982 letter
from G.~P\'olya to A.~Odlyzko \cite{Polya1982letter}, recounting a
Göttingen conversation with E.~Landau around 1912--1914; the first
published mention appears only in Montgomery's 1973 pair-correlation
paper \cite{Montgomery1973}; standard survey and Millennium-problem
accounts include \cite{Bombieri2000,Conrey2003}; and the conjecture
has been taken up repeatedly since
\cite{Berry1985, BerryKeating1999, KeatingSnaith2000, Connes1999}. It
has driven much of the analytic number theory of the last century.
Yet, as Connes' 2026 survey \cite{Connes2026} emphasises, no such
operator has been explicitly constructed for $\zeta(s)$.

The programme implicitly makes a branch choice. For the Selberg zeta
function $Z_\Gamma(s)$ of a compact hyperbolic surface
$\Gamma\backslash\mathbb{H}$, the Laplace-Beltrami operator
$\Delta_{\Gamma\backslash\mathbb{H}}$ realises the Hilbert-P\'olya
structure explicitly: its eigenvalues $\lambda_n = 1/4 + r_n^2$ give
$r_n = \operatorname{Im}(\rho_n)$ for the Selberg zeros $\rho_n =
1/2 + ir_n$, and it is compatible with the canonical trace-formula
transfer algebra generated by the geodesic/isometry flow
\cite{Selberg1956}. For the Riemann zeta function, by contrast,
Results NE-A and NE-B from the recent Even Dominance proof programme
\cite{Geiger2026RHII} rule out the two most natural algebraic
candidates: the prime-shift Fourier multiplier
$M(\xi) = -2\,\Re\zer{\halfline + i\xi}$ fails to be absolutely
totally positive on a set of positive measure (NE-A), and no
non-scalar symmetric matrix simultaneously commutes with all
prime-shift matrices at Galerkin truncation $N\leq 15$ (NE-B, with
full-space extension conjectural \cite{Geiger2026RHII}). And for the
conjectural Prime-Hub graph (Open Problem 5 of
\cite{Geiger2026RHIII}), four obstructions---Weyl-law incompatibility
with the logarithmic density $N(T)\sim (T/2\pi)\log(T/2\pi)$, trace-class
regularity of a graph transfer operator, compatibility with NE-A,
and determinant orientation between $1/\zeta(s)$ and ordinary
Fredholm zeros---keep an explicit Prime-Hub/Hilbert-P\'olya graph construction
open.  After the Spectral Zookeeper \cite{Zookeeper2026}, this
openness is no longer the status of the Riemann $C^{(2)}$/MS2 step
itself; it is the status of the direct OP5 graph interpretation.

\subsection{The trinity of meta-principles}
\label{sec:trinity}

These three canonical facts jointly suggest that the Hilbert-P\'olya
property is not a universal feature of zeta-type families, but a
\emph{local} invariant, present in some sub-families and structurally
absent in others. The central organising principle of this paper is
that this branch-locality is governed by a trinity of meta-principles
that each RH-type programme must address:

\begin{description}
\item[(UCU)] \emph{Universal Convexity Uniqueness.} The admissible
  states of the programme are the minimisers of a strictly convex,
  coercive functional. Strict convexity guarantees a \emph{unique}
  saturation state, which the spectral machinery of the Weil quadratic
  form then realises. This principle organises the uniqueness
  arguments across four otherwise unrelated programmes
  (turbulence~K41, FST dark-energy potential, Riemann Even Dominance,
  Yang--Mills Gibbs measure) under a single head.

\item[(SGE)] \emph{Semigroup-Group Equivalence.} Whether a natural
  self-adjoint operator---the sought-for Hilbert-P\'olya Hamiltonian
  $H_X$---exists is conjecturally governed by a single algebraic
  provenance invariant: whether the transfer structure has locally
  compact group provenance or only semigroup-like provenance. In the
  realised geometric cases, this is the condition under which a Lie
  structure and hence a Casimir operator emerges.

\item[(WF)] \emph{Weil-Form Branch-Transversality.} The Weil
  quadratic form $Q_X$ is well-defined for every zeta-type family
  with an Euler product, a gamma factor, and a functional equation,
  independently of the $\HBL$-status. It provides a uniform
  architecture that yields gap positivity in all algebraic regimes.
  In the Riemann case, the $Q_\zeta$-based Even-Dominance
  companion programme \cite{Geiger2026RHII} is currently framed as a
  conditional reduction and diagnostic route toward the Riemann
  hypothesis that uses no Hilbert-P\'olya input.
\end{description}

In short:

\begin{quote}
\emph{UCU selects the unique saturation state; SGE conjecturally asks when
a natural Hilbert-P\'olya operator exists; Weil-form transversality
builds the bridge across the SGE dichotomy.}
\end{quote}

The remainder of this paper formalises this trinity, tests it on
five zeta-type families, and closes with a speculative reading in
which the three principles mirror the algebraic architecture of the
Bitcoin blockchain protocol (Section~\ref{sec:blockchain}).

\subsection{A field guide to the zeta zoo}
\label{sec:zoo-framing}

Throughout this paper we use the \emph{zeta zoo} as a pedagogical
frame. Each zeta-type family $X$ is an enclosure in the zoo; the
enclosure types are the three possible HP-status values
(\text{YES}, \text{NO}, \text{OPEN}); the zoo-wide feed is the Weil
quadratic form $Q_X$, which every enclosure accepts regardless of
type; the zoo guide is the trinity UCU + SGE + Weil-form
transversality, and the four-class rigor hierarchy
(Section~\ref{sec:rigor}). Model enclosures (Riemann, Selberg,
Prime-Hub) are the main tour stops. Recent additions
(Section~\ref{sec:sge-dedekind-test}: Dedekind
$\mathbb{Q}(\sqrt{-5})$; Section~\ref{sec:sge-ihara-test}: Ihara
Petersen) sit as first non-canonical zoo exhibits. Species that
also inhabit the adjacent FST-physical biotope
\cite{RFEPFramework}---\emph{bridge species} such as the Riemann
zeta itself, the BSD L-function, the Hodge motive, and the Boolean
algebra of P vs NP---are cross-referenced at the relevant entries.

The zoo metaphor is a reading aid, not a theorem scaffold. All
mathematical assertions below are formulated technically; the
metaphor is used solely to orient the reader through the variety
of zeta-type families that the trinity organises.

\subsection{The CRM circle: from cosmos to algebra and back}
\label{sec:crm-circle}

Historically, the Cooperative Renormalization Model (CRM) of
cosmological dark energy was the first instantiation of what later
became Functional Stability Theory \cite{RFEPFramework}. Its
composition law
\[
  x \oplus y = \frac{x + y}{1 + xy}
\]
is the Lorentz-boost addition, so the underlying geometry is
hyperbolic by construction. When the CRM stability argument was
generalised, FST emerged; when FST was abstracted further, the
mathematical trinity of UCU, SGE, and Weil-form transversality
crystallised. The circle closes: from cosmos to physics to
principle to algebra, and the hyperbolic geometry that CRM realises
concretely reappears on the YES-side of the SGE dichotomy as the
$\PSL_2(\mathbb{R})$-Casimir of Selberg zeta
(Section~\ref{sec:selberg}) and, on both sides, as the
branch-transversal Weil-form architecture
(Section~\ref{sec:weil-transversal}). At the end, everything is
algebra---including the geometry.

\subsection{Organization}
\label{sec:intro-organization}

Section~\ref{sec:setup} recalls the definitions of the
Hilbert-P\'olya structure and the Weil quadratic form in the
zeta-family context. Section~\ref{sec:locality} introduces
branch-locality formally and embeds it in the trinity. Sections
\ref{sec:riemann}, \ref{sec:selberg}, and \ref{sec:primehub} recall
the three canonical instances: Riemann (NE-A and NE-B), Selberg
(Casimir operator), and Prime-Hub (Open Problem 5).
Section~\ref{sec:boundary} states and proves the boundary theorem.
Section~\ref{sec:ucu} introduces UCU formally and exhibits its
role as the universal uniqueness statement across four programmes.
Section~\ref{sec:weil-transversal} discusses Weil-form methods as
the branch-transversal complement. Section~\ref{sec:sge-conjecture}
formulates the SGE conjecture and tests it on the Dedekind and
Ihara families. Section~\ref{sec:blockchain} closes with the
blockchain reading, which transports the essence of the trinity
into an informatic register. Section~\ref{sec:outlook} collects
the rigor hierarchy, the consequences for Dirichlet, and the
outlook on adjacent projects.

% ============================================================
\section{Setup}
\label{sec:setup}

\subsection{Zeta-type families}
\label{sec:zeta-family}

\begin{definition}[Zeta-type family]\label{def:zeta-family}
A \emph{zeta-type family} is a tuple $(Z_X, \mathcal{T}_X, \mathcal{E}_X)$
consisting of
\begin{enumerate}
\item a meromorphic function $Z_X(s)$ on $\mathbb{C}$ with at most
  finitely many poles, non-trivial zeros $\{\rho_k = 1/2 + i\gamma_k^X\}$
  on or near a critical line $\{\Re(s) = 1/2\}$, satisfying an Euler
  product on a right half-plane and a functional equation
  $Z_X(s) \leftrightarrow Z_X(1-s)$ up to a gamma factor;
\item a natural family $\mathcal{T}_X$ of bounded operators on a
  Hilbert space associated with the geometry/arithmetic of $X$
  (``transfer operators'');
\item an explicit formula $\mathcal{E}_X$ relating test-function
  sums over zeros to prime/geodesic sums of the family.
\end{enumerate}
\end{definition}

We use standard analytic-number-theory terminology for Euler products,
functional equations and automorphic $L$-functions as in
\cite{IwaniecKowalski}.

\begin{example}[Canonical instances]\label{ex:canonical}
(i) $X = \zeta$: $\mathcal{T}_\zeta = \{\text{prime-shift operators}
\; D(r)\}$ with explicit formula of Weil type.\\
(ii) $X = Z_\Gamma$ for $\Gamma$ co-compact Fuchsian: $\mathcal{T}_\Gamma$
is the canonical trace-formula/wave-kernel algebra induced by the
geodesic flow, with Selberg explicit formula.\\
(iii) $X = G_\mathrm{prime}$: conjectural transfer operator on a
prime-indexed graph; Open Problem 5 of \cite{Geiger2026RHIII}.
\end{example}

\subsection{The Hilbert-P\'olya structure}

\begin{definition}[Branch-local Hilbert-P\'olya]\label{def:hp-bl}
A zeta-type family $(Z_X, \mathcal{T}_X, \mathcal{E}_X)$ is said to
\emph{possess Hilbert-P\'olya structure}---written $\HBL(X) =
\mathrm{YES}$---if there exists a self-adjoint operator
$H_X : \mathcal{D}(H_X) \to \mathcal{H}_X$ on a Hilbert space
$\mathcal{H}_X$ such that:
\begin{enumerate}
\item[(HP1)] $\Spec(H_X) = \{\gamma_k^X : k \in \mathbb{N}\}$, with
  multiplicities.
\item[(HP2)] $H_X$ commutes with every $T \in \mathcal{T}_X$ in the
  sense of commuting spectral projections.
\item[(HP3)] $H_X$ is \emph{provenance-natural}: its admissible
  candidate class and normalisation are fixed from the geometry,
  arithmetic, or operator provenance of $X$ independently of the
  target zero ordinates, and the assignment is functorial under
  provenance-preserving unitary equivalence.
\end{enumerate}
If no such $H_X$ exists, we write $\HBL(X) = \mathrm{NO}$;
if its existence is open, $\HBL(X) = \mathrm{OPEN}$.
\end{definition}

\begin{remark}[Informal intuition]\label{rem:hp-informal}
Hilbert and P\'olya's original suggestion was that the Riemann
hypothesis would follow from exhibiting a natural self-adjoint
operator ``in Nature'' whose eigenvalues coincide with the imaginary
parts of the zeta zeros. The three conditions (HP1)--(HP3) formalise
this informal requirement: spectral identity, commutation with the
natural transfer operators of the family, and a naturality condition
that prevents tautological constructions (e.g.\ the trivial
multiplication operator on $\ell^2$ indexed by the zeros themselves).
\end{remark}

\subsection{The Weil quadratic form as branch-transversal architecture}
\label{sec:weil-qf}

For any zeta-type family $(Z_X, \mathcal{T}_X, \mathcal{E}_X)$ with a
completed functional equation and an explicit formula in the usual
Weil sense, the Weil explicit formula \cite{Weil1952} induces a
quadratic form $Q_X$ on the chosen test-function class of schematic
shape
\begin{equation}
\label{eq:weil-qf-generic}
  Q_X[f] = \sum_k 2\pi |\widehat{f}(\gamma_k^X)|^2
    = W^X_\infty[f] + (\text{prime/geodesic side}).
\end{equation}
After the standard completion conventions have been fixed (gamma
factor, poles, trivial zeros, and exceptional spectral terms),
positivity of $Q_X$ on the relevant test-function class is the
Weil-criterion form of the RH-type statement for $Z_X$, independently
of the existence of any Hilbert-P\'olya operator. This is the
branch-transversal content of the Weil architecture: it applies in
the HP-BL-YES, -NO, and -OPEN cases, but only after those local
normalisations have been specified.

\subsection{Three-component decomposition of RH-type statements}
\label{sec:three-component}

A further structural refinement, following
\cite{MetaGalerkinBridge}, isolates three conceptually distinct
layers inside every RH-type statement for a zeta-type family $X$.
Writing $\mathrm{RH}(X)$ for the positivity of $Q_X$ on the
admissible test-function class, we use the following proof-obligation
route ledger, not a theorem-level equivalence:
\begin{equation}
\label{eq:three-component}
  \mathsf{RHRoute}(X):\qquad
  \underbrace{\mathrm{GBC}(X)}_{\text{Galerkin-Bridge}}
  \;\longrightarrow\;
  \underbrace{C^{(2)}(X)}_{\text{eigenvector realisation}}
  \;\longrightarrow\;
  \underbrace{\mathrm{Limit}_{\mathrm{Hurwitz}}(X)}_{\text{universal limit}}.
\end{equation}
\eqref{eq:three-component} records which three layers every RH-type
programme must supply, and labels them as the targets of distinct
technical sub-programmes. In each concrete family the labels must be
instantiated with their own domain, compactness, convergence and
non-vanishing hypotheses; the precise conventions are fixed in the
individual instances of Sections~\ref{sec:riemann}--\ref{sec:primehub}.

The three components are as follows.

\begin{enumerate}
\item[\textbf{(GBC)}] \emph{Galerkin-Bridge component.} The
algebraic-spectral step: for every cutoff $\lambda > 1$ one
constructs a finite-dimensional Galerkin approximation $Q_X^{(N,\lambda)}$
of $Q_X$ on a logarithmic host space (cosine/sine basis on
$[-\log\lambda,\log\lambda]$ in the arithmetic case), shows
positivity/negativity of the approximation via interval arithmetic
or combinatorial arguments, and transfers the result to the
continuous operator via Mosco convergence
\cite{MetaGalerkinBridge,ConnesConsaniMoscovici2025}. This layer is where the
v2.1 Even-Dominance / Direct-Frontier reduction route
\cite{Geiger2026RHII} lives for $X = \zeta$.

\item[$\boldsymbol{C^{(2)}}$] \emph{Eigenvector realisation
component.} The analytic step: one shows that the Galerkin
eigenvectors $f_\star^{(N)}$ converge, in a sufficiently strong
topology, to a genuine eigenfunction $f_\star$ of the continuous
operator. For the Riemann case this is precisely the Paley--Wiener
control required in Connes' programme
\cite{Connes2026,ConnesConsaniMoscovici2025} and flagged in \cite{Geiger2026RHII}
as the deepest remaining analytic ingredient before the Zookeeper
closure.  In the current CORE status, \cite{Zookeeper2026} closes this
$C^{(2)}$/MS2 step in the CCM/Fourier model and supplies the
microcluster normal form used below.  For Selberg this step is
automatic (compact resolvent of the Laplacian on
$\Gamma\backslash\mathbb{H}$).

\item[\textbf{(Hurwitz)}] \emph{Universal limit component.} The
final functorial step: Hurwitz' theorem on locally uniform limits
of non-vanishing holomorphic functions promotes the family of
finite-level statements to the limit on the critical line. This
component is independent of $X$: it is the same analytic lemma
that enters every $\lambda$-cutoff RH argument, which is why we
list it separately as ``universal''.
\end{enumerate}

\begin{remark}[Where each FST branch specialises]
\label{rem:three-component-specialisation}
The three-component decomposition makes visible \emph{where} each
branch of FST concentrates its technical effort:
\begin{itemize}
\item The \textbf{v2.1 Even Dominance programme}
\cite{Geiger2026RHII} is currently framed as a conditional
reduction/diagnostic route for $X = \zeta$: finite CAP computations
and the GBC layer are evidence and normalisation data, while the public
RH conclusion remains conditional on the odd-sector lower-bound /
AVG-style closure and the associated $C^{(2)}(\zeta)$ bridge.
\item The \textbf{Connes programme}
\cite{Connes2026,ConnesConsaniMoscovici2025} is the natural home of the
$C^{(2)}$ step, formulated semilocally as an analytic condition on
prolate-spheroidal eigenfunction asymptotics; in the Zookeeper paper
this step is closed in the CCM/Fourier branch by the three-lemma
microcluster argument.
\item The \textbf{Obstruction-Classifier programme}
\cite{GeigerObstructionClassifiers} diagnoses on which
(GBC, $C^{(2)}$) pair a given problem is hard: easy GBC +
Zookeeper-closed $C^{(2)}$ (Riemann), twisted-transfer
$C^{(2)}$ (Dirichlet), easy both (Selberg, automorphic
$L$ on $GL_n$), open GBC (Ihara generic, Yang--Lee, $p$-adic
$L$).
\item The \textbf{Hurwitz} layer is invariant across all these
branches; when we ``take the limit'' to the continuous critical
line, the same lemma does the work.
\end{itemize}
In particular, the trinity of Section~\ref{sec:ucu}---UCU + SGE +
Weil-form transversality---acts on the GBC layer (UCU supplies
uniqueness of the Galerkin ground state, SGE classifies the
admissible transfer algebras, Weil-form transversality is the
architecture carrying the Galerkin approximation). $C^{(2)}$ and
Hurwitz are transversal to this trinity; they enter \emph{after}
the trinity has fixed the ground state.
\end{remark}

\begin{remark}[The boundary-theorem view]
\label{rem:boundary-three-component}
Theorem~\ref{thm:boundary} (HP-BL) classifies zeta-type families
by their HP status. The three-component decomposition refines this
by telling us \emph{which component determines the classification}:
\begin{itemize}
\item HP-BL $=$ YES (Selberg, Ihara Petersen, automorphic):
$C^{(2)}$ is automatic (compact resolvent); the HP operator is
the natural Laplacian.
\item HP-BL $=$ NO (Riemann, Dedekind $\mathbb{Q}(\sqrt{-5})$,
Dirichlet): HP-BL failure/no-construction is not identical with an
open $C^{(2)}$ status.  For Riemann the Zookeeper paper closes
$C^{(2)}$/MS2 in the CCM/Fourier branch; for Dirichlet and Dedekind
the corresponding twisted transfer remains to be tested.  The
Weil-form route remains branch-transversal.
\item HP-BL $=$ OPEN (Prime-Hub): the explicit OP5 graph construction
and its boundary-operator interpretation remain open; this is no
longer read as an open Riemann $C^{(2)}$ blocker after Zookeeper.
\end{itemize}
\end{remark}

% ============================================================
\section{Branch-locality as a concept}
\label{sec:locality}

Branch-locality refers to the observation that a structural property
of zeta-type families can hold for one family and fail for a closely
related one. Before we state the boundary theorem, we make the
concept precise.

\subsection{Three modes of \texorpdfstring{$\HBL$}{HP-BL}}

\begin{definition}[HP-BL status trichotomy]\label{def:hpbl-status}
For a zeta-type family $X$ in the sense of Definition~\ref{def:zeta-family},
we say that
\begin{itemize}
\item $\HBL(X) = \mathrm{YES}$ if there exists an operator $H_X$
  satisfying (HP1)--(HP3) of Definition~\ref{def:hp-bl};
\item $\HBL(X) = \mathrm{NO}$ if one can prove, or has established
  under a specific conjecture, that no such $H_X$ exists;
\item $\HBL(X) = \mathrm{OPEN}$ if no candidate has been exhibited
  \emph{and} no exclusion has been proved, possibly subject to
  identified obstructions.
\end{itemize}
\end{definition}

The status $\mathrm{NO}$ may be \emph{unconditional} or
\emph{conditional}. In the present work, the Riemann case will be
$\mathrm{NO}$ conditional on Conjecture~\ref{conj:NE-B-full} (the
extension of NE-B from finite-dimensional Galerkin truncations to the
full space). The conjecture is conjectural but structurally precise: an explicit
open statement whose resolution would either confirm or refute the
NO-classification.

\subsection{Naturality (HP3): the operational test}

The naturality clause (HP3) of Definition~\ref{def:hp-bl} is the
subtle one. Without it, a trivial ``HP operator'' always exists: on
the Hilbert space $\ell^2(\{\gamma_k\})$ one can simply define $H_X$
as multiplication by $\gamma_k$. Such an ad hoc $H_X$ satisfies (HP1)
tautologically but fails to explain anything. The operative safeguard
is therefore not absolute uniqueness inferred from the target
spectrum; it is source-independent provenance and equivariance.

\begin{definition}[Provenance-natural HP operator]\label{def:natural-hp}
Fix a provenance package
\[
 \mathfrak{P}_X=(\mathcal{H}_X,\mathcal{A}_X,\pi_X,
 \mathcal{T}_X,\mathcal{C}_X,\nu_X),
\]
where the symmetry/operator algebra $\mathcal{A}_X$, its representation
$\pi_X$, the transfer algebra $\mathcal{T}_X$, the admissible
self-adjoint candidate class $\mathcal{C}_X$, and the sign, scale, and
additive normalisation $\nu_X$ are specified from the geometry,
arithmetic, or operator provenance of $X$ without using the target
zero ordinates. A candidate $H_X\in\mathcal{C}_X$ is
\emph{provenance-natural} if:
\begin{enumerate}
\item[(N1)] the package is fixed before testing (HP1), so the zero list
  is an output test rather than construction data;
\item[(N2)] every spectral projection of $H_X$ commutes with
  $\pi_X(\mathcal{A}_X)$ and $\mathcal{T}_X$;
\item[(N3)] a unitary equivalence $U$ of provenance packages transports
  the candidate functorially, $H'_X=UH_XU^*$;
\item[(N4)] the normalisation $\nu_X$ is fixed before (HP1), and any
  additional uniqueness assertion is supported by a separate rigidity
  certificate relative to $\mathcal{C}_X$ rather than inferred from
  the prescribed spectrum.
\end{enumerate}
\end{definition}

\begin{lemma}[Unitary provenance invariance]\label{lem:hp3-unitary}
If $U$ is a unitary equivalence of provenance packages and $H_X$ obeys
(N1)--(N4), then $UH_XU^*$ is self-adjoint, has the same spectrum with
the same multiplicities, and obeys the transported strong commutation
relations. Hence the HP-BL label is invariant inside a fixed
provenance-equivalence class.
\end{lemma}

\begin{proof}
Unitary conjugation preserves self-adjointness and spectrum. For every
Borel set $B$, the spectral theorem gives
$E_{UH_XU^*}(B)=U E_{H_X}(B)U^*$. Intertwining the provenance and
transfer representations therefore transports every commutation
relation, while the definition of a provenance equivalence transports
$\mathcal{C}_X$ and $\nu_X$.
\end{proof}

The zero-indexed multiplication operator fails (N1) unless an
independent geometric or arithmetic construction supplies both its
basis and its operator. In the Selberg row, the ambient group is
$G=\PSL_2(\mathbb{R})=\operatorname{Isom}^+(\mathbb{H})$, whereas
$\Gamma\backslash\mathbb{H}$ is the quotient surface, not a symmetry
group. Its hyperbolic metric canonically supplies the Laplace-Beltrami
operator and bounded functional calculus; the curvature $-1$ and the
explicit $1/4$ shift fix the normalisation
\cite{Selberg1956,Marklof2008}. In the Riemann row,
Theorem~\ref{thm:NE-B} establishes only that no non-scalar symmetric
operator commutes with every $D_N(r)$ at the tested truncations in the
prime-shift provenance class. The Prime-Hub row would have to provide
its own independent graded, groupoid, or dynamical provenance, which
remains part of Open Problem~5.

\subsection{Why not a uniform theorem across all zeta-type families?}

A plausible naive hope is that $\HBL$ would be \emph{uniformly}
$\mathrm{YES}$ across zeta-type families: every $L$-function would
carry its own natural $H_X$. The boundary theorem below is a precise
quantification of the failure of this hope: even the three canonical
instances $\{\zeta, Z_\Gamma, G_\mathrm{prime}\}$ already realise all
three values $\{\mathrm{NO}, \mathrm{YES}, \mathrm{OPEN}\}$. Any
uniform claim across zeta-type families is therefore false.

A subtler hope is that $\HBL = \mathrm{YES}$ would be \emph{generic},
with the Riemann case somehow exceptional. Theorem~\ref{thm:NE-B}
disallows even this reading on the tested Galerkin truncations: the
absence of a common symmetry among prime-shift operators is an
explicit structural feature of the multiplicative prime structure,
not a pathology. Conversely, the presence of a Lie-group-invariant
Laplacian in Selberg's setting is a consequence of the specific
geometry of $\Gamma\backslash\mathbb{H}$, not of zeta-analyticity per
se. The branch-locality thesis is that $\HBL$-status is determined
by the \emph{algebraic substrate} of the family, independently of
its zeta-analytic features.

\begin{remark}[Relative status]\label{rem:relative-hbl}
The notation $\HBL(X)$ is always relative to the declared transfer
presentation $\mathcal{T}_X$ in Definition~\ref{def:zeta-family}.
When ambiguity matters one should read it as $\HBL_{\mathcal{T}}(X)$.
Thus the Riemann row below is a statement about the classical
prime-shift transfer presentation. Connes' adelic programme, discussed
below, is a different transfer presentation and remains a possible
open route to an adelic $\HBL_{\mathcal{T}}(\zeta)=\mathrm{YES}$.
\end{remark}

\begin{remark}[Strict and resonance YES]\label{rem:yes-types}
We distinguish two positive statuses. $\mathrm{YES}_{\mathrm{sa}}$
denotes the strict self-adjoint, point-spectrum Hilbert--P\'olya
realisation of Definition~\ref{def:hp-bl}. $\mathrm{YES}_{\mathrm{res}}$
denotes the non-compact flow-zeta analogue in which a natural
self-adjoint operator controls a scattering/resolvent problem and the
zeros are resonances rather than ordinary eigenvalues. The latter is
useful for the zeta-zoo taxonomy, but it is not a literal instance of
(HP1) without this resonance qualification.
\end{remark}

\subsection{A three-class taxonomy of the zeta zoo}
\label{sec:three-class-taxonomy}

The branch-locality trichotomy
$\HBL \in \{\mathrm{YES}, \mathrm{NO}, \mathrm{OPEN}\}$ of
Definition~\ref{def:hpbl-status} classifies families by the
\emph{existence} of a natural Hilbert-P\'olya operator. A second,
largely orthogonal classification captures the \emph{type of spectrum}
that such an operator, if it exists, admits. Three classes emerge:

\begin{center}
\small
\setlength{\tabcolsep}{3pt}
\renewcommand{\arraystretch}{1.18}
\begin{tabular}{@{}p{0.19\linewidth} p{0.18\linewidth} p{0.12\linewidth} p{0.18\linewidth} p{0.17\linewidth}@{}}
\toprule
\textbf{Class} & \textbf{Transfer algebra} & \textbf{Typical HP} & \textbf{Spectrum} & \textbf{Zoo examples} \\
\midrule
Arithmetic zetas & Prime/prime-ideal semigroup & NO or SGE-predicted NO & Discrete (zeros) & Riemann, Dirichlet, Dedekind \\
Geometric zetas (compact) & Discrete $\Gamma$ on compact quotient & YES$_{\mathrm{sa}}$ & Discrete (Laplace) & Selberg, Ihara \\
Flow zetas (non-compact) & Continuous Lie group on non-compact space & YES$_{\mathrm{res}}{}^*$ & Continuous + resonances & CRM, Ruelle-Anosov, BH-QNM \\
\bottomrule
\end{tabular}
\end{center}
${}^*$ realised in the resonance sense: the HP operator (Casimir of
the acting Lie group) has absolutely continuous spectrum, and the
zeta zeros correspond to scattering resonances in the lower half
plane (cf.~\cite{DyatlovZworski2019}).

The three classes are orthogonal to the $\HBL$-trichotomy: both
Selberg and CRM sit on positive sides of the SGE dichotomy, but
Selberg is $\mathrm{YES}_{\mathrm{sa}}$ while CRM is only
$\mathrm{YES}_{\mathrm{res}}$. Riemann and
Dedekind share the arithmetic class but remain separate from the
geometric and flow classes entirely. We will exhibit one canonical
member of each class in
Sections~\ref{sec:riemann}--\ref{sec:crm-flow}, and give the
complete boundary theorem in Section~\ref{sec:boundary}.

\subsection{Relation to Connes' adelic framework}

In Connes' programme \cite{Connes1999, Connes2026}, the relevant
``space'' on which a putative $H_\zeta$ acts is an adelic quotient
$\mathbb{A}_K / K^\times$, and the natural symmetry is the
$\mathrm{GL}_1$-action of idele classes. Our formulation
Definition~\ref{def:natural-hp} does not exclude this architecture:
if Connes' programme succeeds, his operator would satisfy (N1)--(N3)
with $G_X = \mathrm{GL}_1(\mathbb{A}_K/K^\times)$, and would
constitute an explicit $\HBL = \mathrm{YES}$ certificate for
$\zeta(s)$ in the adelic Hilbert space. That possibility is currently
open: it would coexist with Theorem~\ref{thm:NE-B} at finite
dimensions because the adelic natural $G_X$ is not visible in the
Galerkin truncation. Our boundary theorem therefore concerns
Hilbert-P\'olya structure with respect to the \emph{prime-shift}
transfer operators $\mathcal{T}_\zeta$, which is the relevant
structure for the Weil quadratic form and the Even Dominance proof programme
of~\cite{Geiger2026RHII}.

The sections that follow exhibit one canonical instance of each
$\HBL$ mode.

% ============================================================
\section{Instance A: Riemann (NE-A and NE-B)}
\label{sec:riemann}

\subsection{Non-existence of absolute total positivity (NE-A)}
\label{sec:ne-a}

We quote the following theorem from \cite{Geiger2026RHII}:

\begin{theorem}[NE-A, {\cite[Thm.~NE-A]{Geiger2026RHII}}]\label{thm:NE-A}
The Fourier multiplier
$M(\xi) = -2\,\Re[\zeta'/\zeta(\halfline + i\xi)]$
of the prime-shift operator takes negative values on a set of
positive Lebesgue measure in every interval $[-L, L]$ with $L > \gamma_1$
(the first positive zeta ordinate). Consequently, the prime-shift
operator is not absolutely totally positive ($\mathrm{PF}_\infty$).
\end{theorem}

\begin{remark}\label{rem:ne-a-empirical}
Empirically (cf.\ \cite[\S4.2]{Geiger2026RHII}), sampling $M(\xi)$
on a uniform grid of $10^4$ points in $\xi \in [0,100]$ with
$\varepsilon$-neighbourhoods of the first $29$ ordinates excluded
yields approximately $49\%$ of points with $M(\xi) < 0$. This
illustrates the theorem; it does not replace the structural argument.
\end{remark}

\subsection{No universal commuting operator (NE-B)}
\label{sec:ne-b}

\begin{theorem}[NE-B finite certificate, {\cite[Thm.~NE-B]{Geiger2026RHII}}]\label{thm:NE-B}
For every registered Galerkin truncation $N \leq 15$ and the certified
dense shift samples $r_1, \ldots, r_k \in (0,2)$ reported in
\cite{Geiger2026RHII}, the computed common symmetric centraliser of
the prime-shift matrices $D_N(r_j)$ is one-dimensional. Equivalently,
within those finite certified truncations, every simultaneously
commuting symmetric matrix is a scalar multiple of the identity.
\end{theorem}

\begin{conjecture}[NE-B-full, {\cite[Conj.~\ref{thm:NE-B}-full]{Geiger2026RHII}}]\label{conj:NE-B-full}
Theorem~\ref{thm:NE-B} extends to $N = \infty$: no non-trivial
universal commuting operator exists for the Riemann prime-shift family.
\end{conjecture}

\subsection{Consequence for \texorpdfstring{$\HBL(\zeta)$}{HP-BL(zeta)}}

Theorems~\ref{thm:NE-A} and~\ref{thm:NE-B} jointly exclude the two
most natural algebraic routes to (HP1)--(HP3) for the classical
prime-shift presentation of the Riemann zeta family. No
absolutely totally positive prime shift operator can serve as $H_\zeta$
(Route A), and no universal commuting self-adjoint symmetric operator
exists at any tested Galerkin truncation (Route B). Under
Conjecture~\ref{conj:NE-B-full}, $\HBL(\zeta) = \mathrm{NO}$.

\begin{remark}[Paper II's conditional RH route does not require HP structure]
\label{rem:riemann-without-hp}
The Even-Dominance route to the Riemann hypothesis in
\cite{Geiger2026RHII} is currently treated as a conditional
reduction/diagnostic route: it proceeds by prime-weighted summation
(Shift Parity Lemma) plus frontier-dominance arguments using the Weil
quadratic form, but the finite CAP evidence is not by itself an
unconditional RH proof. Neither step requires (HP1)--(HP3). This is
consistent with the branch-locality thesis: even if
$\HBL(\zeta) = \mathrm{NO}$ (conditional on
Conjecture~\ref{conj:NE-B-full}), the Weil-form architecture can remain
the relevant non-HP route, subject to its own odd-sector/AVG closure.
\end{remark}

\begin{remark}[Documentation status of the Direct-Frontier route]
\label{rem:riemann-critique-review}
We cite \cite{Geiger2026RHII} as the companion programme that documents
the Even-Dominance / Direct-Frontier reduction route. A systematic
critique of the trilogy
\cite{Geiger2026RHII,Geiger2026RHIII}---covering (K1) the global
extrapolation from the $\lambda = 100$ CAP baseline, (K2) the
transparency of the PNT/Mertens constants
$C_{\mathrm{shift}}, C_{\mathrm{norm}}$, (K3) the Galerkin convergence
to the infinite-dimensional Weil operator, (K4) the
implication-versus-equivalence structure of the reduction chain A1--A8,
(K5) the precise characterisation of the form domain of $Q_\zeta$, and
(K6) overall reviewer ergonomics---is treated here as identifying the
remaining guardrails of that route. In the current status those
guardrails are not cosmetic: the public RH conclusion is conditional
until the odd-sector lower-bound / AVG-style closure and the relevant
Galerkin-to-Weil bridge are supplied. This qualification does not affect
the boundary theorem or the SGE classification used in the present
paper; it only prevents the companion programme from being used here as
an unconditional RH proof.
\end{remark}

% ============================================================
\section{Instance B: Selberg (Casimir operator)}
\label{sec:selberg}

\subsection{The Laplace-Beltrami as explicit Hilbert-P\'olya operator}

Let $\Gamma \subset \PSL_2(\mathbb{R})$ be a co-compact Fuchsian group
and $\mathcal{F} = \Gamma\backslash\mathbb{H}$ the associated compact
hyperbolic surface. The Selberg zeta function $Z_\Gamma(s)$ has
non-trivial zeros $\rho_n = \halfline + ir_n$ related to the spectrum
of the Laplace-Beltrami operator $\Delta_\mathcal{F}$ on
$L^2(\mathcal{F})$ via
\begin{equation}
\label{eq:selberg-spectral}
  \Delta_\mathcal{F}\,\phi_n = \lambda_n\phi_n, \qquad
  \lambda_n = \tfrac{1}{4} + r_n^2,
\end{equation}
with finitely many exceptions from small eigenvalues
\cite{Selberg1956,Marklof2008}.

For this row, $\mathcal{T}_\Gamma$ denotes the trace-formula /
wave-kernel algebra generated by the functional calculus of
$\sqrt{\Delta_\mathcal{F}-1/4}$; closed geodesics enter through the
Selberg trace formula. We do not assert commutation with arbitrary
Ruelle, symbolic, or weighted transfer operators on anisotropic spaces.

\begin{theorem}[Selberg]\label{thm:selberg-hp}
After separating the standard finite exceptional sector
$\{\lambda_n<1/4\}$, the operator
$H_{Z_\Gamma}:=\sqrt{\Delta_\mathcal{F}-1/4}$ on the spectral subspace
$\lambda_n\geq 1/4$ satisfies (HP1)--(HP3) of
Definition~\ref{def:hp-bl} for the principal critical-line part of
the zeta-type family
$(Z_\Gamma, \mathcal{T}_\Gamma, \mathcal{E}_\Gamma)$. The exceptional
real zeros are a finite Selberg correction, not part of the
Hilbert-P\'olya ordinate spectrum. In this standard qualified sense,
$\HBL(Z_\Gamma)=\mathrm{YES}$.
\end{theorem}

\begin{proof}[Sketch]
(HP1) follows from~\eqref{eq:selberg-spectral} on the subspace
$\lambda_n\geq 1/4$, where $r_n$ is real. For (HP2),
$\mathcal{T}_\Gamma$ is the declared bounded functional-calculus /
wave-kernel algebra of $\Delta_\mathcal{F}$, so its elements strongly
commute with the spectral projections of $H_{Z_\Gamma}$. For (HP3),
the curvature-$-1$ hyperbolic metric fixes $\Delta_\mathcal{F}$
independently of the Selberg zeros; the Laplacian on $\mathbb{H}$
commutes with isometry pullbacks and therefore descends to
$\Gamma$-invariant functions on $\Gamma\backslash\mathbb{H}$
\cite{Marklof2008}. This is the provenance-natural construction of
Definition~\ref{def:natural-hp}; no uniqueness claim is inferred from
the zero spectrum. Since $\Gamma\backslash\mathbb{H}$ is compact in
this instance, $\Delta_\mathcal{F}$ has discrete spectrum; the finitely
many small eigenvalues are treated as the usual exceptional Selberg
zeros rather than by an absolutely continuous spectral part.
\end{proof}

\subsection{Selberg trace formula: the explicit mechanism}
\label{sec:selberg-trace}

The Selberg trace formula \cite{Selberg1956} schematically relates
spectral and geometric data:
\begin{equation}
\label{eq:selberg-trace}
  \sum_n \widehat{h}(r_n)
  = \frac{\operatorname{vol}(\mathcal{F})}{4\pi} \int h(r)\,r\tanh(\pi r)\,dr
  + \sum_{\{\gamma\}} \frac{\ell(\gamma_0)}{2\sinh(\ell(\gamma)/2)}\,
      g(\ell(\gamma)),
\end{equation}
where $r_n$ are the spectral parameters of $\Delta_\mathcal{F}$, the
second sum runs over primitive closed geodesics $\{\gamma_0\}$ and
their iterates, and $h, g$ are a standard test-function pair related
by Fourier transform. The geodesic side of \eqref{eq:selberg-trace}
is the \emph{multiplicative-geometric analogue} of the prime side of
the Riemann--Weil explicit formula.

The metric invariance of $\Delta_\mathcal{F}$ is the structural reason
that the spectral side has a canonical functional calculus. The
closed geodesics on the geometric side of
\eqref{eq:selberg-trace} are indexed by hyperbolic conjugacy classes
and their lengths; they are trace data, not a separately asserted
family of commuting operators $T_{\ell_\gamma}$
\cite{Selberg1956,Marklof2008}. Equivalently, the Casimir provenance
of the ambient $\PSL_2(\mathbb{R})$ induces the normalised Laplacian,
while the theorem uses only its bounded functional-calculus algebra.
This is the concrete verification of (N1)--(N4) in the safe scope of
Definition~\ref{def:natural-hp}.

\subsection{What distinguishes Selberg from Riemann}
\label{sec:selberg-vs-riemann}

The Riemann zeta family admits a formal analogue of
\eqref{eq:selberg-trace} in the shape of the Weil explicit formula:
\begin{equation}
\label{eq:weil-riemann}
  \sum_k h(\gamma_k^\zeta)
  = W_\infty[h]
  - 2\sum_{p,m} \frac{\log p}{p^{m/2}}\,\widehat{h}(m\log p),
\end{equation}
with archimedean weight $W_\infty$ encoding the gamma factor. The
prime side of \eqref{eq:weil-riemann} is \emph{multiplicative}: it is
indexed by the semigroup of prime powers, not by orbits of a Lie
group action. There is no natural continuous group whose discrete
orbits reproduce $\{p^m\}$, and Theorem~\ref{thm:NE-B} confirms
this structurally: no non-scalar symmetric operator commutes with
the prime-shift operators at the tested Galerkin truncations.

The precise structural difference between
\eqref{eq:selberg-trace} and \eqref{eq:weil-riemann} is therefore:
\begin{center}
\begin{tabular}{@{}lll@{}}
\toprule
Feature & Selberg (HP-BL YES) & Riemann (HP-BL NO${}^*$) \\
\midrule
Indexing set & closed geodesics $\{\gamma\}$ &
  prime powers $\{p^m\}$ \\
Group action & $\PSL_2(\mathbb{R})$ on $\mathbb{H}$ & none (NE-B) \\
Commuting Hamiltonian & $\Delta_\mathcal{F}$ = Casimir & none at $N \leq 15$ \\
Length structure & continuous $\ell(\gamma) \in \mathbb{R}_{>0}$ &
  discrete $\log p \in \mathbb{R}_{>0}$ \\
\bottomrule
\end{tabular}
\end{center}

The multiplicative arithmetic of the primes is \emph{not the orbit
structure of any natural continuous symmetry group}. This is the
content of branch-locality: the algebraic substrate of the
transfer-operator family decides the HP-BL status, independently
of the zeta-analytic similarity between \eqref{eq:selberg-trace}
and \eqref{eq:weil-riemann}.

% ============================================================
\section{Instance C: Prime-Hub (Open Problem 5)}
\label{sec:primehub}

Open Problem~5 of \cite{Geiger2026RHIII}, hereafter OP5, calls for an
explicit topologically mixing directed graph
$G_\mathrm{prime} = (V, E)$ (or equivalently a flow or dynamical system)
satisfying five conditions:
\begin{enumerate}
\item[(OP5.1)] Primitive periodic orbits $\{\gamma_p\}$ are in
  bijection with the primes.
\item[(OP5.2)] $\ell(\gamma_p) = \log p$ for every prime $p$.
\item[(OP5.3)] The associated transfer operator $\mathcal{L}_s$ has
  Fredholm determinant $\det(I - \mathcal{L}_s)^{\mathrm{vac}} =
  1/\zeta(s)$.
\item[(OP5.4)] The spectral radius of $\mathcal{L}_s$ is less than $1$
  for $\Re(s) > 1/2$.
\item[(OP5.5)] Eigenvalue $1$ occurs exactly at the non-trivial zeros
  of $\zeta(s)$.
\end{enumerate}
No construction simultaneously satisfying (OP5.1)--(OP5.5) is known
since the informal Hilbert-P\'olya suggestion of 1914
\cite{Polya1982letter}. The programme is reviewed and surveyed in
\cite{Connes2026}; Connes' \emph{Letter Through Time} (\emph{ibid.},
final sections) proposes a trace-formula-based convergence strategy
from finite to infinite Euler products, but a constructive
$G_\mathrm{prime}$ is not exhibited.

\subsection{Obstruction C1: Weyl-law mismatch}
\label{sec:obs-C1}

\begin{proposition}[Weyl-law mismatch]\label{prop:weyl-obstruction}
Let $H$ be the Laplace-Beltrami operator of a smooth compact
$d$-dimensional Riemannian manifold. Then its eigenvalue counting
function satisfies Weyl's law
\begin{equation}
\label{eq:weyl-classical}
  N_H(T) := \#\{\lambda \in \Spec(H) : \lambda \leq T^2\}
  \;\sim\; \frac{\omega_d \operatorname{vol}(M)}{(2\pi)^d}\,T^d
  \qquad (T \to \infty),
\end{equation}
with $\omega_d$ the volume of the unit ball in $\mathbb{R}^d$. The
Riemann-zero counting function is
\begin{equation}
\label{eq:riemann-N}
  N_\zeta(T) := \#\{\gamma_k^\zeta : 0 < \gamma_k^\zeta \leq T\}
  =
  \frac{T}{2\pi}\log\!\frac{T}{2\pi e} + O(\log T)
\end{equation}
(classical Riemann-von Mangoldt, cf.\ \cite[\S9.4]{Titchmarsh1986}).
The asymptotic rates \eqref{eq:weyl-classical} and
\eqref{eq:riemann-N} are incompatible for every integer $d \geq 1$.
In particular, no classical geometric Laplacian on a smooth compact
manifold can realise $\Spec(H_\zeta)$.
\end{proposition}

\begin{proof}
The power-law \eqref{eq:weyl-classical} contrasts with the
logarithmic density \eqref{eq:riemann-N}: for any $d \geq 1$, the
ratio $N_\zeta(T) / N_H(T) \to 0$ (if $d \geq 2$) or $\to \infty$ as
$T \to \infty$ (if $d = 1$), so no choice of $d$ and $\operatorname{vol}(M)$
makes \eqref{eq:weyl-classical} match \eqref{eq:riemann-N}.
Consequently, any candidate $H_\zeta$ with
$\Spec(H_\zeta) = \{\gamma_k^\zeta\}$ cannot be the Laplacian of a
smooth compact Riemannian manifold.
\end{proof}

\begin{remark}\label{rem:weyl-workarounds}
Proposition~\ref{prop:weyl-obstruction} rules out smooth compact
manifold Laplacians but not all possibilities: non-compact spectra,
hyperbolic-like volume growth (Berry's chaotic quantisation
\cite{Berry1985, BerryKeating1999}), or adelic spaces
\cite{Connes1999, Connes2026} remain candidate architectures. The
obstruction is therefore a restriction on the \emph{type} of
geometric substrate.
\end{remark}

\subsection{Obstruction C2: Trace-class regularity}
\label{sec:obs-C2}

\begin{proposition}[Trace-class regularity]\label{prop:trace-class}
The Fredholm determinant $\det(I - \mathcal{L}_s)^{\mathrm{vac}}$ in
requirement (OP5.3) is well-defined only if $\mathcal{L}_s$ is
trace-class on its ambient Hilbert space. For a graph transfer
operator whose primitive orbits have logarithmic lengths
$\ell(\gamma_p) = \log p$, the summability condition
\begin{equation}
\label{eq:trace-class-condition}
  \sum_{p,m} \frac{1}{p^{m/2}} \cdot |\lambda_p^m|
  \;<\; \infty
\end{equation}
(with $\lambda_p$ the eigenvalue of $\mathcal{L}_s$ along the primitive
cycle $\gamma_p$) must hold uniformly on the regime where (OP5.4)
demands spectral radius less than $1$. No explicit construction
simultaneously satisfying (OP5.1)--(OP5.5) and
\eqref{eq:trace-class-condition} has been exhibited; this is known in
the Geiger tree as obstruction~K4 of \cite{Geiger2026RHI}.
\end{proposition}

\begin{remark}\label{rem:trace-class-connection}
The difficulty behind Proposition~\ref{prop:trace-class} is that
(OP5.4) and \eqref{eq:trace-class-condition} pull in opposite
directions: (OP5.4) constrains the eigenvalues $\lambda_p$ from above
by $1$, while \eqref{eq:trace-class-condition} constrains their
\emph{sum} to be absolutely convergent. For a ``generic'' graph
with infinitely many primitive orbits, natural operator models give
either violation of (OP5.4) (eigenvalues pile up near $1$) or
violation of \eqref{eq:trace-class-condition} (slow decay). Threading
this needle is the substantive content of OP5.
\end{remark}

\subsection{Obstruction C3: NE-A compatibility}
\label{sec:obs-C3}

\begin{proposition}[NE-A compatibility]\label{prop:ne-a-compat}
Any ordinary positive trace-formula Hilbert-P\'olya candidate
$H_\mathrm{prime}$ for $\zeta$, in the prime-shift presentation,
implicitly determines the Fourier multiplier
$M(\xi) = -2\,\Re[\zeta'/\zeta(1/2 + i\xi)]$ of the prime-shift
operator via the trace formula of $(H_\mathrm{prime}, \mathcal{L}_s)$.
By Theorem~\ref{thm:NE-A}, $M(\xi) < 0$ on a set of positive
Lebesgue measure. Consequently, $H_\mathrm{prime}$ cannot arise as
a non-negative self-adjoint operator with non-negative Fourier
symbol on smooth geometries; it must carry intrinsic sign structure
in its spectral transform.
\end{proposition}

\begin{proof}[Sketch]
Such a positive trace-formula candidate must reproduce the Weil explicit
formula (Section~\ref{sec:weil-qf}) via its trace. The kernel of the trace
formula is precisely $M(\xi)$. Standard Laplacians on smooth geometry
produce non-negative Fourier multipliers; the sign of $M(\xi)$ on a
positive-measure set therefore rules out such Laplacians.
Construction must use structures with inherent sign change (e.g.\
pseudo-differential, or noncommutative adelic).
\end{proof}

\begin{remark}\label{rem:ne-a-connection}
Obstructions C1 (smooth geometry ruled out) and C3 (non-negative
symbols ruled out) are independent but mutually reinforcing: C1
rules out manifold architecture by density, C3 rules out manifold
architecture by symbol sign. Even if one workaround (e.g.\
hyperbolic chaotic quantisation) circumvents C1, it must still
engineer a symbol satisfying C3.
\end{remark}

\subsection{Obstruction C4: determinant orientation}
\label{sec:obs-C4}

\begin{proposition}[Fredholm orientation]\label{prop:fredholm-orientation}
Let $U$ be a neighbourhood of a non-trivial zero $\rho$ of $\zeta$,
and let $s\mapsto\mathcal{L}_s$ be a holomorphic trace-class family on
$U$. If
\[
  D(s) := \det(I-\mathcal{L}_s)
\]
is an ordinary Fredholm determinant and (OP5.3) is read after
meromorphic continuation as $D(s)=1/\zeta(s)$ near $\rho$, then
$\rho$ cannot be an ordinary eigenvalue-one event of $\mathcal{L}_s$.
Indeed, an eigenvalue-one event forces $D(\rho)=0$, whereas
$1/\zeta(s)$ has a pole at $\rho$.
\end{proposition}

\begin{proof}
For a holomorphic trace-class family, the Fredholm determinant is
holomorphic. If $1$ enters the spectrum of $\mathcal{L}_\rho$, then
$I-\mathcal{L}_\rho$ is not invertible and the ordinary Fredholm
determinant vanishes at $\rho$. On the other hand, if $\rho$ is a zero
of $\zeta$ of order $m\geq 1$, then $1/\zeta(s)$ has a pole of order
$m$ at $\rho$. A holomorphic determinant zero and a meromorphic pole
cannot be the same local event. Therefore the two OP5 signals cannot
live in one ungraded ordinary Fredholm determinant.
\end{proof}

\begin{remark}[Graded Prime-Hub reading]\label{rem:graded-primehub}
The obstruction does not rule out OP5; it fixes its semantics. The
symbol ``vac'' in (OP5.3) must encode a quotient, regularised, or
graded determinant, or OP5 must be split into two channels:
\[
  D_E(s)=\prod_p(1-p^{-s})=\frac{1}{\zeta(s)},\qquad
  D_Z(s)=\frac{\xi(s)}{\xi(0)},\qquad
  G(s)=\frac{\xi(s)}{\xi(0)\zeta(s)}
\]
with
\[
  D_Z(s)=\frac{G(s)}{D_E(s)}.
\]
The Euler channel $D_E$ carries prime orbits and M\"obius parity; the
spectral channel $D_Z$ carries the non-trivial zero signal; and
$G$ contains the completion factors. Equivalently, one may seek a
$\mathbb{Z}_2$-graded transfer complex with
\[
  \operatorname{Sdet}(I-\mathcal{L}_s)
  =\frac{\det(I-\mathcal{L}_s^+)}
          {\det(I-\mathcal{L}_s^-)}
  =\frac{1}{\zeta(s)},
\]
so that eigenvalue-one events at zeta zeros appear in the denominator
sector. The elementary finite model is the M\"obius-graded prime
operator $A_s e_p=p^{-s}e_p$, for which
\[
  \det(I-A_s)=\operatorname{Str}_{\Lambda^\ast\ell^2(P)}(\Lambda A_s)
  =\sum_{n\ \mathrm{squarefree}}\mu(n)n^{-s}
  =\prod_p(1-p^{-s})
\]
for $\Re(s)>1$ in the infinite prime set. This is recorded in the
public computational note \cite{MobiusGradedPrimeHub}. The current numerical
follow-up \cite{PrimeWeilBridgeDefect} shows why a non-tautological
Prime--Weil defect must replace any direct insertion of a zero list.
\end{remark}

\subsection{Consequence}

The combination of Propositions~\ref{prop:weyl-obstruction},
\ref{prop:trace-class}, \ref{prop:ne-a-compat}, and
\ref{prop:fredholm-orientation} identifies four
structural obstructions that any explicit Prime-Hub construction
must navigate. None of the four is individually a proof of
impossibility; together, they delimit the narrow parameter space in
which a valid $G_\mathrm{prime}$ could exist. Consequently,
$\HBL(G_\mathrm{prime}) = \mathrm{OPEN}$, with obstructions C1--C4
as explicit guides for any future construction. Connes'
\emph{Letter Through Time} \cite{Connes2026} exemplifies the kind of
non-smooth, noncommutative architecture that could in principle
navigate all four; whether the programme closes remains open.

% ============================================================
\section{Instance D: CRM as Flow-Zeta (FST-Cosmology bridge)}
\label{sec:crm-flow}

Instances A, B, C exhibit the arithmetic and geometric-compact
classes of the three-class taxonomy of Section~\ref{sec:three-class-taxonomy}.
The flow-zeta class is populated by a concrete object from the
physical side of FST: the Cooperative Renormalization Model (CRM)
of cosmological dark energy \cite{RFEPFramework}. Its mathematical
structure is developed in full detail in \cite{CRMFlowZetaNote}; we
summarise the key points here.

\subsection{CRM as a one-parameter Lie group}

The CRM composition law governing dark-energy evolution is
\begin{equation}
\label{eq:crm-composition}
  x \oplus y = \frac{x+y}{1+xy}, \qquad x, y \in (-1, 1).
\end{equation}
Setting $x = \tanh u$, $y = \tanh v$ gives $x \oplus y = \tanh(u+v)$,
so $(x, \oplus)$ is isomorphic to the additive group $(\mathbb{R},
+)$. The composition is therefore the Lorentz-boost addition, and
the relevant transfer algebra is a \emph{one-parameter Lie group}
realised as a subgroup of $\PSL_2(\mathbb{R})$ acting on the
hyperbolic plane $\mathbb{H}$.

\begin{proposition}[CRM transfer algebra]\label{prop:crm-transfer}
For the CRM family $X = \mathrm{CRM}$, the transfer algebra is
$(I_{\mathrm{CRM}}, \cdot) = (\mathbb{R}, +)$, a locally compact
abelian Lie group. Consequently, under Conjecture~\ref{conj:sge},
$\HBL(\mathrm{CRM}) = \mathrm{YES}_{\mathrm{res}}$ in the
non-compact resonance sense of Remark~\ref{rem:yes-types}.
\end{proposition}

\subsection{The natural Hilbert-P\'olya operator and its spectrum}

The Casimir of $\mathfrak{sl}_2(\mathbb{R})$ restricted to the
boost-generating component is, up to equivalence, the
Laplace-Beltrami operator $\Delta_{\mathbb{H}}$ on the hyperbolic
plane \cite{DyatlovZworski2019}. In non-compact underlying geometry,
this is not a strict (HP1) point-spectrum realisation. Rather,
$\Delta_{\mathbb{H}}$ supplies the natural self-adjoint provenance
and commutation structure for a resonance theory: its absolutely
continuous spectrum is $[1/4, \infty)$, and it commutes with the full
$\PSL_2(\mathbb{R})$-action, hence with the CRM flow in particular.

The \emph{zeros} of $\zeta_{\mathrm{CRM}}$ are, in this setting,
\emph{scattering resonances}: poles of the meromorphic continuation
of the resolvent $(\Delta_{\mathbb{H}} - s(s-1))^{-1}$ through the
essential spectrum into the lower half plane. On asymptotically
hyperbolic backgrounds (e.g.\ de~Sitter, Schwarzschild--de~Sitter)
these resonances are precisely the \emph{quasinormal modes} of the
relevant wave equation \cite{DyatlovZworski2019}. This realises
$\HBL = \mathrm{YES}_{\mathrm{res}}$ in the resonance sense of
Section~\ref{sec:three-class-taxonomy}.

\subsection{The Ruelle flow-zeta of CRM}

Following \cite{DyatlovZworski2019}, the appropriate analytic object
is the Ruelle-Dyatlov-Zworski zeta
\begin{equation}
\label{eq:crm-flow-zeta}
  \zeta_{\mathrm{CRM}}(s)
  = \prod_{\gamma \in \mathcal{P}_{\mathrm{CRM}}}
    \frac{1}{1 - e^{-s\,\ell(\gamma)}},
\end{equation}
where $\mathcal{P}_{\mathrm{CRM}}$ indexes the primitive periodic
orbits of the CRM flow on an asymptotically hyperbolic background,
and $\ell(\gamma)$ is the orbit length in the Lorentz-boost metric.
The zeros of $\zeta_{\mathrm{CRM}}(s)$ are the resonances described
above.

\subsection{CRM as a bridge species}

CRM occupies a distinctive position in the zeta zoo: it is
simultaneously a flow-zeta entry in the mathematical zoo (this
section) and a Pattern~A instantiation in the physical RFEP
ecosystem \cite{RFEPFramework} (see also the narrative closure in
Section~\ref{sec:crm-circle}). Together with Riemann (historical
origin), BSD, Hodge, and P~vs~NP, CRM forms the fifth
\emph{bridge species} of the programme---and the first to open the
flow-zeta column on both sides of the FST twins.

\subsection{The CRM saturation theorem as a classification
statement}
\label{sec:crm-saturation-classification}

A structural remark is in order. The CRM~V saturation theorem
\cite{CRMV} asserts that any dynamical system on $(-1,1)$ satisfying
four axioms---existence of a scalar relaxation variable, existence of
a dissipative term, existence of a stable saturation value, and
monotone smooth dynamics---necessarily realises the $\tanh$ saturation
profile. Read as a characterisation of the flow-zeta class, this is
not merely a statement about one cosmological model; it is a
\emph{classification statement} for the class of IR-asymptotics
compatible with functional stability on the boost interval.

This is structurally analogous, on the flow side, to the Hurwitz/
compactness characterisation used in the arithmetic branch: just as
the Weil quadratic form detects stability across branches, CRM~V
delineates the admissible asymptotic class on the flow-noncompact
branch. In particular, several candidate UV completions may
realise distinct selections from this saturation class; the
classification does not depend on any individual realisation.

This observation matters for two reasons. First, it frames CRM not
as a bespoke cosmological proposal but as an IR-sector calibration of a
branch-local structural pattern. Second, it suggests the audit test that any UV
theory whose IR limit passes the four CRM~V axioms (e.g.\ a
GUT-level geometric model, an RG-flow construction, or an
open-system holographic construction) lands in the same saturation
class---modulo quantitative parameter matching. Making this
UV/IR compatibility test operational is an open programme within
FST-Cosmology.

% ============================================================
\section{The Boundary Theorem}
\label{sec:boundary}

\begin{theorem}[Branch-Locality of Hilbert-P\'olya: Conditional Core]\label{thm:boundary}
Conditional on Conjecture~\ref{conj:NE-B-full} for the Riemann row and
relative to the prime-shift transfer presentation, the $\HBL$-property
is not constant across zeta-type families. When combined with the
three-class taxonomy of Section~\ref{sec:three-class-taxonomy}, the zoo
exhibits three theorem-relevant canonical rows and several calibration
rows. The table is an evidence ledger: only the Riemann, Selberg and
Prime-Hub rows enter the theorem-facing classification; the CRM,
Dedekind and Ihara rows calibrate the taxonomy and carry different
evidence classes.

\begin{center}
\renewcommand{\arraystretch}{1.15}
\begin{tabular}{@{}p{0.19\linewidth}p{0.16\linewidth}p{0.21\linewidth}p{0.33\linewidth}@{}}
\toprule
Family & Status & Scope / class & Evidence class \\
\midrule
Riemann $\zeta(s)$ & NO${}_{\mathrm{prime}}{}^*$ & arithmetic prime-shift &
  Thms.~\ref{thm:NE-A}, \ref{thm:NE-B}; Conj.~\ref{conj:NE-B-full} \\
Selberg $Z_\Gamma(s)$ & YES$_{\mathrm{sa}}$ & geometric-compact & theorem-level: Thm.~\ref{thm:selberg-hp} \\
Prime-Hub $G_\mathrm{prime}$ & OPEN & --- &
  open with obstructions C1--C4: Props.~\ref{prop:weyl-obstruction}--\ref{prop:fredholm-orientation} \\
CRM (Dark Energy) & YES$_{\mathrm{res}}$ & flow-noncompact & resonance analogue: Prop.~\ref{prop:crm-transfer}, \cite{DyatlovZworski2019} \\
Dedekind $\mathbb{Q}(\sqrt{-5})$ & SGE-predicted NO & arithmetic finite test & finite NE-B evidence ($N\leq 10$): \S\ref{sec:sge-dedekind-test}, \cite{DedekindNEBTest} \\
Ihara $\zeta_{\text{Petersen}}$ & finite YES proxy & finite graph & finite certificate/proxy: \S\ref{sec:sge-ihara-test}, \cite{IharaPetersenTest} \\
\bottomrule
\end{tabular}
\end{center}
${}^*$ conditional on Conjecture~\ref{conj:NE-B-full} and relative to
the prime-shift transfer presentation.
\end{theorem}

\begin{proof}
The theorem-facing classification already follows from the canonical
Riemann, Selberg, and Prime-Hub rows of
Sections~\ref{sec:riemann}--\ref{sec:primehub}: a conditional
prime-shift NO, a strict self-adjoint YES, and an OPEN obstruction
case are distinct. The CRM row is a resonance analogue supplied by
Section~\ref{sec:crm-flow}; the Dedekind and Ihara rows are finite
SGE stress tests in Sections~\ref{sec:sge-dedekind-test} and
\ref{sec:sge-ihara-test}. They provide additional calibration and
do not enter as proof assumptions for the theorem.
\end{proof}

\begin{corollary}[Implication for the Riemann hypothesis]
\label{cor:rh-consequence}
Any unconditional proof of the Riemann hypothesis that relies on
(HP1)--(HP3) for the classical prime-shift presentation of $\zeta$
must first refute
Conjecture~\ref{conj:NE-B-full}. Alternatively, Weil-form methods
(Section~\ref{sec:weil-transversal}) provide a separate
conditional/diagnostic companion route that does not require the
Hilbert-P\'olya structure; this route is documented in
\cite{Geiger2026RHII} but is not used here as an unconditional proof.
\end{corollary}

% ============================================================
\section{Universal Convexity Uniqueness (UCU)}
\label{sec:ucu}

This section isolates the uniqueness principle that operates, in
precisely the same abstract form, in all programmes of the
functional-stability landscape considered in this paper. We call
it \emph{Universal Convexity Uniqueness} (UCU). It is the
mathematical lemma underlying what is called \emph{Pattern~A} in
the physical formulation (RFEP; \cite{RFEPFramework}) and its
dynamical partner \emph{Dissipative Selection} (DS1--DS3). UCU is
a single technical statement; its universality lies in how many
seemingly unrelated objects realise it.

\subsection{The UCU lemma}
\label{sec:ucu-lemma}

\begin{lemma}[UCU]\label{lem:ucu}
Let $(\mathcal{H}, \langle\cdot,\cdot\rangle)$ be a real Hilbert
space, let $\mathcal{A} \subset \mathcal{H}$ be a closed convex
admissible class, and let $J : \mathcal{A} \to \mathbb{R}$ satisfy
\begin{enumerate}
\item[(C)] \emph{Strict convexity:} for all $u,v \in \mathcal{A}$
with $u \neq v$ and all $t \in (0,1)$,
\[
  J(tu + (1-t)v) < t\,J(u) + (1-t)\,J(v);
\]
\item[(L)] \emph{Strong lower semicontinuity:} if $u_n \to u$
strongly in $\mathcal{H}$ with $u_n, u \in \mathcal{A}$, then
$J(u) \leq \liminf_{n\to\infty} J(u_n)$;
\item[(K)] \emph{Coercivity on $\mathcal{A}$:} $J(u) \to +\infty$
whenever $\|u\| \to \infty$ within $\mathcal{A}$.
\end{enumerate}
Then $J$ has a unique minimiser $u_\star \in \mathcal{A}$.
\end{lemma}

\begin{proof}
Let $\{u_n\}$ be a minimising sequence. By (K) it is bounded, so it
admits a weakly convergent subsequence $u_{n_k} \rightharpoonup u_\star$
with $u_\star \in \mathcal{A}$ by weak closedness of closed convex
sets. Convexity (C) plus strong lower semicontinuity (L) together
imply weak lower semicontinuity of $J$ (Mazur's lemma; see
\cite[Prop.~II.1.2]{EkelandTemam1999}), so
$J(u_\star) \leq \liminf J(u_{n_k}) = \inf_{\mathcal{A}} J$, and
$u_\star$ is a minimiser. Strict convexity (C) excludes a second
minimiser: if $v_\star \neq u_\star$ also realised the minimum,
the midpoint $\tfrac{1}{2}(u_\star + v_\star) \in \mathcal{A}$
would strictly decrease $J$, contradicting minimality.
\end{proof}

The content of Lemma~\ref{lem:ucu} is classical, but the scope of
its role in the present programme is not: it is the \emph{same}
lemma, with different $(\mathcal{H}, \mathcal{A}, J)$, that yields
the unique saturation state in each of the five realisations
listed below.

\subsection{The parallel realisations}
\label{sec:ucu-instances}

Table~\ref{tab:ucu-instances} exhibits the five instances.
Each column shows which object plays the role of the Hilbert space
$\mathcal{H}$, the admissible class $\mathcal{A}$, and the
functional $J$. The last row records the unique minimiser $u_\star$
and its physical or arithmetic interpretation.
The table is a template ledger, not a claim that every row has the
same proof status. Some rows are theorem-level convex minimisation
statements; others are programme-level normal forms whose strict
convexity/coercivity hypotheses remain part of the referenced
companion project.

\renewcommand{\arraystretch}{1.25}
\begin{table}[htbp]
\centering
\scriptsize
\setlength{\tabcolsep}{3pt}
\begin{tabular}{|p{0.16\linewidth}|p{0.40\linewidth}|p{0.34\linewidth}|}
\hline
\textbf{Programme} & \textbf{Functional $J$ on class $\mathcal{A}$} &
\textbf{Unique minimiser $u_\star$} \\ \hline

\emph{Riemann Even Dominance} \cite{Geiger2026RHII}
& $J(f) = \langle f, Q_\zeta f \rangle / \|f\|^2$ on
$\mathcal{A} = L^2_\mathrm{even}([-L,L]) \setminus \{0\}$;
convexity on even test functions reflects the positive-definiteness
of the Weil quadratic form after Direct Frontier Dominance.
& The ground state $f_\star$ realising the spectral gap;
its existence and uniqueness are exactly the
even--odd dominance statement. \\ \hline

\emph{Selberg zeta (Theorem~\ref{thm:selberg-hp})}
& $J(\phi) = \int_{\Gamma\backslash\mathbb{H}} |\nabla \phi|^2$
on $\mathcal{A} = \{\phi \in H^1_0(\mathcal{F}) : \|\phi\|_{L^2}=1\}$;
the Rayleigh quotient is strictly convex on the sphere after
mod-out of the $S^1$-phase.
& The first Laplace--Beltrami eigenfunction $\phi_1$, whose
Casimir eigenvalue supplies the spectral Hilbert--P\'olya content. \\ \hline

\emph{Turbulence / Kolmogorov K41}
& $J(u) = \int |\nabla u|^2 - \epsilon \int u$ on
$\mathcal{A} = $ admissible divergence-free velocity fields with
fixed energy dissipation $\epsilon$; strict convexity holds in the
inertial range via the exponent-$2/3$ scaling.
& The unique scaling profile producing the $k^{-5/3}$ energy
cascade \cite{FSTTurbulence}. \\ \hline

\emph{Yang--Mills Gibbs measure}
& $J(A) = S_{\mathrm{YM}}(A) = \tfrac{1}{4g^2}\int F \wedge {\star F}$
on $\mathcal{A} = \mathcal{A}_{\mathrm{flat}} / \mathcal{G}$, the
moduli space of gauge-equivalent connections; strict convexity on
the quotient follows from the Coulomb gauge condition.
& The vacuum sector achieving the mass gap
\cite{FSTYangMills}. \\ \hline

\emph{CRM (flow-zeta, \S\ref{sec:crm-flow})}
& $J(\mathrm{rapidity}) = V_{\mathrm{eff}}(\mathrm{rapidity})$, an
effective dark-energy potential on $(-1,1)$ with boundary
coercivity at $|x| \to 1$ built in by the Lorentz-boost composition
$\oplus$; strict convexity is the structural input of the
saturation axioms \cite{CRMV}.
& The unique saturated boost state, producing the observed $\tanh$
profile in $w(z)$. \\ \hline

\end{tabular}
\caption{UCU instances as a template ledger. Same lemma, different
$(\mathcal{H},\mathcal{A},J)$; row status is inherited from the
referenced companion programme. Where strict convexity and coercivity
are established on the stated admissible class, UCU yields the unique
saturation state.}
\label{tab:ucu-instances}
\end{table}
\renewcommand{\arraystretch}{1.0}

\subsection{UCU and Pattern~A / RFEP: the translation}
\label{sec:ucu-pattern-a}

The physical face of UCU in the Renormalized Free Energy Principle
framework \cite{RFEPFramework} is \emph{Pattern~A}: functional
positivity under a gauge constraint. The following is the formal
translation template:

\begin{itemize}
\item \textbf{Pattern~A (PA1).} A target functional $J$ (free
energy, Weil form, Rayleigh quotient, K41 action, Yang--Mills
action) is positive-definite on a gauge-restricted state space.\\
$\longleftrightarrow$ \emph{UCU hypothesis (C).}

\item \textbf{Pattern~A (PA2).} The functional is bounded below
and grows at infinity of the admissible class.\\
$\longleftrightarrow$ \emph{UCU hypothesis (K).}

\item \textbf{Pattern~A (PA3).} A unique stable saturation exists
and is attained.\\
$\longleftrightarrow$ \emph{UCU conclusion.}
\end{itemize}

\begin{remark}[UCU vs.\ DS1--DS3]\label{rem:ucu-vs-ds}
The dissipative selection principle DS1--DS3 of
\cite{RFEPFramework} is the \emph{dynamical} partner of UCU: UCU
identifies the unique minimiser; DS1--DS3 identifies the unique
trajectory to it. In all five instances of
Table~\ref{tab:ucu-instances} a time-dependent or flow-parametrised
version of $J$ admits the DS1--DS3 structure, but the minimiser
itself is already fixed by UCU alone. DS1--DS3 are therefore
logically downstream of UCU.
\end{remark}

\subsection{UCU and the trinity}
\label{sec:ucu-trinity}

With UCU in place, the trinity structure of this paper reads:
\[
  \underbrace{\text{UCU}}_{\text{uniqueness of saturation}}
  \;+\;
  \underbrace{\text{SGE}}_{\text{HP-status classification}}
  \;+\;
  \underbrace{\text{Weil-form transversality}}_{\text{branch-independent positivity tool}}
\]

UCU provides the \emph{why is there a unique answer} layer; SGE
(Section~\ref{sec:sge-conjecture}) provides the \emph{which branch does the
family belong to} layer; the Weil-form architecture
(\S\ref{sec:weil-transversal}) provides the \emph{how to compute
positivity in each branch} layer. The three layers are
independent: UCU does not predict YES or NO on the HP spectrum;
SGE does not predict spectral gaps; Weil-form transversality does
not supply uniqueness by itself. They compose to give the full
programme.

\begin{remark}[Scope of UCU]\label{rem:ucu-scope}
UCU is \emph{necessary but not sufficient} for Hilbert--P\'olya
realisation. A unique minimiser in the sense of
Lemma~\ref{lem:ucu} does not automatically provide a Hermitian
operator whose spectrum encodes the zeros of $Z_X$; it provides a
ground state. The passage from ground state to spectrum requires
the SGE classification (to identify whether one has a group or
only a semigroup of transfer operators) and the Weil-form
architecture (to recover the explicit formula positivity). UCU
is the common denominator, not the proof of Hilbert--P\'olya in
any single instance.
\end{remark}

% ============================================================
\section{Weil-form methods as branch-transversal complement}
\label{sec:weil-transversal}

We now argue that the Weil quadratic form $Q_X$ of
Section~\ref{sec:weil-qf} applies uniformly in all three branch
regimes.

\subsection{Form-domain anchoring}
\label{sec:weil-form-domain}

For each zeta-type family $X$, the Weil quadratic form $Q_X$ is
anchored by the quadruple
\begin{equation}\label{eq:weil-quadruple}
  \bigl(Q_X,\; \operatorname{Dom}(Q_X),\; P_X,\;
        \operatorname{Dom}(Q_X)^{\pm}\bigr),
\end{equation}
where $\operatorname{Dom}(Q_X) \subset L^2(\mathbb{R})$ is a
closed subspace on which $Q_X$ is well-defined and lower-bounded
(the form domain), $P_X$ is the parity operator
$f(t) \mapsto f(-t)$, and
$\operatorname{Dom}(Q_X)^{\pm}
  = \operatorname{Dom}(Q_X) \cap L^2_{\mathrm{even/odd}}$
are the parity-invariant subspaces. For the arithmetic instances
of the present paper (Sections~\ref{sec:riemann}--\ref{sec:primehub})
the form domain is taken as the Friedrichs closure of
$C_c^{\infty}([-\log\lambda,\log\lambda])$ with respect to the
graph norm of $Q_X$, and the cosine/sine bases of
Section~\ref{sec:weil-qf} realise a density system in
$\operatorname{Dom}(Q_X)^{\pm}$. The parity-preservation
$P_X \operatorname{Dom}(Q_X) = \operatorname{Dom}(Q_X)$ is a
consequence of the parity-invariance of the archimedean kernel
$e^{u/2}/|2\sinh u|$ (the absolute value is taken; the Weil
formula uses the principal-value convention at $u = 0$, cf.\
\cite{Weil1952}) and of the symmetric prime-shift operator family
$S_{m\log p} + S_{-m\log p}$ entering the Weil form.

\begin{remark}[Scope of this anchoring]\label{rem:form-domain-scope}
The explicit form-domain lemma underlying the Riemann instance
(restricting $Q_\zeta$ to an $H^{1/2}$-closure on which the
archimedean kernel is form-bounded, and verifying
parity-preservation on the full form domain) is stated as a
minor-revision item in \cite{RHCritiqueReview2026} (K5). The
branch-transversal content of the present paper is independent
of that refinement: the Shift Parity Lemma and the Direct
Frontier Dominance proof operate on finite-dimensional Galerkin
blocks, and their transfer to the continuous operator uses
Cauchy-interlacing plus Schur-complement bounds (see
Section~\ref{sec:three-component}, $C^{(2)}$ component), not
form-domain regularity per se.
\end{remark}

\subsection{Positivity at HP-BL-YES}

For $X = Z_\Gamma$ (Selberg, Theorem~\ref{thm:selberg-hp}), the
positivity of $Q_{Z_\Gamma}$ on the relevant test-function class is
a consequence of the self-adjointness of $\Delta_\mathcal{F}$; the
Weil-form positivity can be derived from the spectral decomposition
of the Laplacian. This is a consistency check: both architectures
give the same answer.

\subsection{Positivity at HP-BL-NO}

For $X = \zeta$ (Riemann, Theorem~\ref{thm:boundary} row 1), the
positivity of $Q_\zeta$ is pursued in the companion programme
\cite{Geiger2026RHII} via the Direct-Frontier / Even-Dominance route
(Shift Parity Lemma + Mertens theorem + CAP-certified base case).
The current public status is conditional, but the route uses no
Hilbert-P\'olya input. This is the strict branch-transversality content.

\subsection{Positivity at HP-BL-OPEN}

For $X = G_\mathrm{prime}$, the Weil-form positivity question splits
into two steps. First, the explicit formula $\mathcal{E}_\mathrm{prime}$
for Prime-Hub must coincide with the Riemann--Weil explicit formula
\eqref{eq:weil-riemann} on the prime/geodesic side, since
requirements (OP5.1)--(OP5.3) of Section~\ref{sec:primehub} prescribe
that the Fredholm determinant of $\mathcal{L}_s$ equals $1/\zeta(s)$.
By Proposition~\ref{prop:fredholm-orientation}, this determinant
identity must be interpreted through a quotient, regularised, or
graded channel if OP5.5 is to remain meaningful. Once that channel
semantics is fixed, $Q_{G_\mathrm{prime}} = Q_\zeta$ as a quadratic form on
the test-function class, and its positivity reduces to that of
Riemann (the HP-BL-NO regime). Second, \emph{provided} an explicit
Prime-Hub construction is eventually produced, the Weil-form
positivity is automatic by the branch-transversal argument of
\cite{Geiger2026RHII} on the common prime side. This is the strongest
possible expression of branch-transversality: the positivity of the
quadratic form is invariant under the $\HBL$-status transition.

Adjacent evidence for a related HP-BL-OPEN family: Session-8
numerical work on Dirichlet $L$-functions
\cite{PaleyWienerDirichlet} shows that attempts to build $H_\chi$
from phenomenological kernels (Gaussian or Hadamard-weighted pole)
fail systematically, while Weil-form positivity remains accessible.
This is a second empirical instance where HP-BL-status is
effectively $\mathrm{NO}$ or $\mathrm{OPEN}$ and Weil-form methods
continue to apply.

\subsection{Summary}

The Weil quadratic form is branch-transversal: its existence and
positivity-content does not depend on the HP-BL status of the family.

% ============================================================
\section{The Semigroup-Group Equivalence (SGE)}
\label{sec:sge-conjecture}

\subsection{A unifying conjecture}
\label{sec:sge-conjecture-statement}

The Boundary Theorem~\ref{thm:boundary} exhibits three qualitatively
distinct values $\{\mathrm{YES}, \mathrm{NO}, \mathrm{OPEN}\}$ of
$\HBL$ on three canonical instances. A natural question is whether
this trichotomy is \emph{fundamental} or merely a surface feature of
a deeper single criterion. We formulate the following unifying
conjecture.

Let $(Z_X, \mathcal{T}_X, \mathcal{E}_X)$ be a zeta-type family in
the sense of Definition~\ref{def:zeta-family}. The transfer family
$\mathcal{T}_X$ is typically organised by a discrete set $I_X$ of
``prime objects'' (rational primes, closed geodesics, primitive
cycles), each carrying a canonical length $\ell(\iota)$. The associated
transfer provenance supplies either a semigroup-like object
(multiplication of prime powers, an ideal monoid, primitive-cycle
concatenation) or a group action with a unitary representation and
trace formula. In the Selberg row, the group side is not a free
composition law on closed geodesics; closed geodesics enter as
trace-formula data, while the Casimir comes from the
$\PSL_2(\mathbb{R})$-invariant geometry.

\begin{conjecture}[Semigroup-Group Equivalence, SGE]
\label{conj:sge}
For the calibrated class of zeta-type families considered in this paper,
\[
  \HBL(X) = \mathrm{YES}
  \quad\Longleftrightarrow\quad
  (I_X, \cdot)\ \text{has locally compact group provenance}.
\]
\end{conjecture}

The backward direction (group $\Rightarrow$ YES) is classical only in
the realised trace-formula cases: a natural group action together with
the appropriate unitary representation and spectral trace formula
produces a Casimir-type operator satisfying (HP1)--(HP3). The SGE
conjecture is the stronger claim that this model-case mechanism
captures exactly the branch-local Hilbert-P\'olya boundary.

\paragraph{Waterline.}
Conjecture~\ref{conj:sge} is a classification conjecture, not a
theorem of this paper. The theorem-level content is the conditional
non-constancy statement of Theorem~\ref{thm:boundary}; CRM, Dedekind,
Ihara-Petersen, and Waisen-type transfers calibrate the taxonomy but do
not upgrade SGE to a proved criterion.

The forward direction (semigroup $\Rightarrow$ NO) is the substantive
content of Conjecture~\ref{conj:sge} and is consistent with every
instance tested to date:

\begin{center}
\begin{tabular}{@{}lcc@{}}
\toprule
Family & Transfer algebra $(I_X, \cdot)$ & Predicted $\HBL$ \\
\midrule
Riemann $\zeta(s)$ & $\{p^m\}$ (commutative semigroup) & NO$_{\mathrm{prime}}^{\mathrm{cond}}$ (Thm.~\ref{thm:boundary}) \\
Dirichlet $L(s,\chi)$ & $\{p^m : p \nmid q\}$ (sub-semigroup) & SGE-predicted NO; no simple HP in \cite{PaleyWienerDirichlet} \\
Dedekind $\zeta_K(s)$ & prime ideals $\{\mathfrak{p}\}$ (ideal monoid) & SGE-predicted NO; finite evidence \\
Selberg $Z_\Gamma(s)$ & group/Casimir trace formula & YES (Thm.~\ref{thm:selberg-hp}) \\
Ihara $\zeta_G(u)$ & closed paths in undirected $G$ & finite YES proxy (Petersen) \\
Automorphic $L(s,\pi)$ & Hecke algebra on $\Gamma\backslash G$ & candidate YES (trace-formula dependent) \\
Prime-Hub $G_\mathrm{prime}$ & primitive cycles (free semigroup) & SGE-predicted NO; current status OPEN \\
\bottomrule
\end{tabular}
\end{center}

If Conjecture~\ref{conj:sge} holds, the apparent trichotomy of
Theorem~\ref{thm:boundary} collapses to a single algebraic criterion.
The OPEN status of Prime-Hub, in this reading, is not a third
fundamental mode but a state of partial knowledge: the explicit
construction of $(I_\mathrm{prime}, \cdot)$ would deterministically
decide between YES and NO via Conjecture~\ref{conj:sge}. A separate
working note \cite{SGENote} develops the precise formulation and
collects falsification criteria.

\subsection[First empirical test of Conjecture: Dedekind zeta of Q(sqrt(-5))]{First empirical test of Conjecture~\ref{conj:sge}:
Dedekind zeta of \texorpdfstring{$\mathbb{Q}(\sqrt{-5})$}{Q(sqrt(-5))}}
\label{sec:sge-dedekind-test}

As a direct probe of Conjecture~\ref{conj:sge} outside the three
canonical instances of Theorem~\ref{thm:boundary}, we performed the
following test \cite{DedekindNEBTest}. For the imaginary quadratic
field $K = \mathbb{Q}(\sqrt{-5})$ (discriminant $-20$, class number
$2$), the Dedekind zeta $\zeta_K$ has a transfer algebra indexed by
prime ideals $\mathfrak{p}$ of $\mathcal{O}_K = \mathbb{Z}[\sqrt{-5}]$,
which is a semigroup under ideal multiplication. Conjecture~\ref{conj:sge}
therefore predicts $\HBL(\zeta_K) = \mathrm{NO}$.

To test this, we enumerated the $92$ prime ideals of $K$ with norm
at most $500$, classified into $2$ ramified (above rational primes
$2, 5$), $86$ split (two ideals for each rational $p \equiv 1, 3, 7,
9 \pmod{20}$), and $4$ inert (norm $p^2$ for $p \equiv 11, 13, 17,
19 \pmod{20}$). The corresponding normalised shifts $r_\mathfrak{p}
= 2\log N\mathfrak{p}/L$ with $L = 10$ yielded $49$ distinct values
in $(0, 2)$---fewer than the $95$ obtained for the rational primes
of $\mathbb{Q}$ up to the same norm bound, owing to the split-prime
doubling collapsing pairs of ideals to a single shift value. For each
shift we constructed the NE-B matrix $D_N(r)$ of
Definition~\ref{thm:NE-B} via numerical quadrature, and computed
the symmetric common centraliser dimension by solving the linear
system $[M, D_N(r_\mathfrak{p})] = 0$ on the $N(N+1)/2$ symmetric
degrees of freedom.

\begin{center}
\begin{tabular}{@{}c c c c@{}}
\toprule
$N$ & $\mathbb{Q}$ (control) & $\mathbb{Q}(\sqrt{-5})$ & random shifts \\
\midrule
$5$ & $\dim(Z) = 1$ & $\dim(Z) = 1$ & $\dim(Z) = 1$ \\
$7$ & $\dim(Z) = 1$ & $\dim(Z) = 1$ & $\dim(Z) = 1$ \\
$10$ & $\dim(Z) = 1$ & $\dim(Z) = 1$ & $\dim(Z) = 1$ \\
\bottomrule
\end{tabular}
\end{center}

All three columns agree: only the scalar centraliser survives at
every tested $N \in \{5, 7, 10\}$. This constitutes finite
controlled evidence for Conjecture~\ref{conj:sge} outside the
three canonical instances of Theorem~\ref{thm:boundary}: despite the
qualitatively distinct prime decomposition pattern of $K$ and the
roughly two-fold lower shift density relative to $\mathbb{Q}$, the
Dedekind transfer algebra of $\mathbb{Q}(\sqrt{-5})$ admits the same
NE-B behaviour as the rational prime-shift family. The observation
supports the algebraic (rather than density-based) nature of NE-B
and, transitively, of the $\HBL = \mathrm{NO}$ prediction.

\subsection{Complementary empirical test of
Conjecture~\ref{conj:sge}: Ihara zeta of the Petersen graph}
\label{sec:sge-ihara-test}

On the $\HBL = \mathrm{YES}$ side of the SGE conjecture, a direct
probe is provided by the Ihara zeta function $\zeta_G(u)$ of an
undirected finite graph $G$. The transfer algebra in this case is the
path groupoid modulo backtracking, whose associated structure is the
fundamental group $\pi_1(G)$---a free group $F_r$ with $r = m - n + 1$
generators. Since $F_r$ is a group (not merely a semigroup),
Conjecture~\ref{conj:sge} predicts a YES-side mechanism for graph
zeta functions. The finite test below only certifies the Petersen
instance as a controlled proxy; it is not a theorem for every connected
undirected graph.

We tested this prediction on the Petersen graph (3-regular,
$n = 10$, $m = 15$, a Ramanujan graph with $\mathrm{Aut}(G) = S_5$).
Three independent verifications were carried out \cite{IharaPetersenTest}:
\begin{enumerate}
\item[(V1)] \textbf{Bass-Hashimoto identity.} The classical identity
  $\det(I - uB) = (1 - u^2)^{m-n}\det(I - Au + Qu^2)$, with $B$ the
  Hashimoto non-backtracking operator on directed edges and
  $Q = D - I$, was verified at the test point $u = 0.3 + 0.1 i$ to
  $15$ significant digits.
\item[(V2)] \textbf{Ramanujan RH-analog.} $18$ of the $19$
  non-trivial zeros of $\zeta_G(u)^{-1}$ lie exactly on the critical
  circle $|u| = 1/\sqrt{q-1} = 1/\sqrt{2}$. The single exception at
  $|u| = 1/2$ corresponds to the trivial adjacency eigenvalue
  $\lambda_A = 3$.
\item[(V3)] \textbf{Finite YES-side certificate/proxy.} The graph Laplacian
  $\Delta_G = D - A$ is self-adjoint with spectrum
  $\{0, 2^{(5)}, 5^{(4)}\}$. Its eigenvalues relate to the Ihara
  zeros via the quadratic pencil $2u^2 - \lambda_A u + 1 = 0$. Every
  graph automorphism $\sigma \in S_5$ induces a permutation matrix
  $P_\sigma$ that commutes with both $\Delta_G$ and the edge-lift
  $P_\sigma^{\mathrm{edge}}$ commutes with $B$. The centraliser of
  $B$ in $\mathrm{End}(\mathbb{C}^{2m})$ therefore contains the
  $S_5$-representation algebra, whose non-scalar elements provide a
  finite YES-side certificate/proxy for
  $\HBL(\zeta_{\mathrm{Petersen}})$.
\end{enumerate}

In combination with the Dedekind test of
Section~\ref{sec:sge-dedekind-test}, Conjecture~\ref{conj:sge} has
now been finite-stress-tested with controls on \emph{both} sides of
the semigroup-group dichotomy: the semigroup side
($\mathbb{Q}(\sqrt{-5})$, $\dim(Z) = 1$) at $N \leq 10$, and the
group side (Petersen, non-scalar commuting operators constructed
explicitly). Five qualitatively distinct zeta-type families now
support Conjecture~\ref{conj:sge} consistently (three canonical
instances plus these two new tests), and no counter-example has been
produced. This remains finite evidence, not a proof. Elevation to a
formally proved theorem would require a structural argument (rather
than further case studies) and is identified as the primary next
step.

\begin{remark}[Discriminating control experiment]
\label{rem:sge-control}
A natural worry about the Dedekind test of
Section~\ref{sec:sge-dedekind-test} is that $\dim(Z) = 1$ might be
the \emph{generic} centraliser dimension of arbitrary matrix
families rather than a structural signature of SGE-NO. A dedicated
control experiment \cite{SGEControlExperiment2026} addresses this
by running the same centraliser-dimension routine on explicit
SGE-YES families (cyclic group $\mathbb{Z}/N$ regular
representation; elementary abelian $(\mathbb{Z}/2)^k$) alongside
a random-symmetric baseline. The outcome is unambiguous:
\begin{itemize}
\item $\mathbb{Z}/N$: $\dim(Z_{\text{sym}}) = 7, 10, 14, 17, 22$
  at $N = 5, 7, 10, 12, 15$ (grows linearly with $N$).
\item $(\mathbb{Z}/2)^k$: $\dim(Z_{\text{full}}) = 2^k$ exactly
  (equal to the group order) for $k = 2, 3, 4$.
\item Random symmetric baseline: $\dim(Z_{\text{sym}}) = 1$ for
  all tested $N$.
\end{itemize}
The test apparatus therefore discriminates cleanly between SGE-YES
(centraliser dimension scales with $N$) and SGE-NO / null
(dimension one). The Dedekind result $\dim(Z) = 1$ is thus
structurally informative, not a generic artefact. This closes a
documentation gap identified in the 7-phase review
\cite{RHCritiqueReview2026} of the present manuscript.
\end{remark}

Connes' adelic programme \cite{Connes1999, Connes2026} fits this
picture naturally: by replacing the semigroup $\{p^m\}$ with the
idele class group $\mathbb{A}_K/K^\times$, it \emph{transports} the
Riemann family from the semigroup column into the group column of
the table. On the SGE reading, Connes' programme would upgrade
$\HBL(\zeta)$ from NO (in the classical prime-shift transfer
algebra) to YES (in the adelic transfer algebra), without
contradicting Theorem~\ref{thm:boundary}.

% ============================================================
\section{The Blockchain Reading: a speculative bridge}
\label{sec:blockchain}

\emph{This section is marked as a \textbf{speculative bridge} in the
four-class rigor hierarchy of Section~\ref{sec:rigor}: it transports
the mathematical content of the trinity into an informatic register
through a correspondence with the Bitcoin blockchain protocol. No
mathematical assertion of this section is new; the contribution is a
consolidated reading that makes the architecture of the trinity
accessible to readers familiar with cryptographic protocols.}

\subsection{Two algebraic worlds in one protocol}
\label{sec:blockchain-two-worlds}

The Bitcoin protocol simultaneously employs two algebraic structures
which are, from the perspective of Conjecture~\ref{conj:sge}, of
opposite type:

\begin{enumerate}
\item The \textbf{hash function} (SHA-256) generates an image set
  with a \emph{semigroup} structure under concatenation of data
  blocks. The hash is deliberately one-way---it has no natural
  inverse operation---so the indexing set of a hash-chain is a
  semigroup, not a group.
\item The \textbf{digital signature} (ECDSA on the elliptic curve
  $\mathrm{secp256k1}$) operates on an \emph{abelian group}: the
  $E(\mathbb{F}_p)$-points form a group under chord-and-tangent
  addition, with explicit inverses and a well-defined Casimir-like
  invariant (the curve equation $y^2 = x^3 + 7$).
\end{enumerate}

In the language of Section~\ref{sec:sge-conjecture}, SHA-256 lives
on the semigroup side of the SGE dichotomy, whereas ECDSA lives on
the group side. The Bitcoin protocol is thus a \emph{physical
realisation} of the algebraic dichotomy that Conjecture~\ref{conj:sge}
identifies at the level of zeta-type families, with the hash side
mirroring the Riemann prime-power semigroup $\{p^m\}$ and the
signature side mirroring the Selberg geodesic group structure
$\Gamma \subset \PSL_2(\mathbb{R})$.

\subsection{The correspondence table}
\label{sec:blockchain-correspondence}

\begin{center}
\renewcommand{\arraystretch}{1.25}
\begin{tabular}{@{}p{0.38\linewidth} p{0.52\linewidth}@{}}
\toprule
\textbf{Tree-theoretic structure} & \textbf{Bitcoin-protocol analog} \\
\midrule
Prime-power semigroup
  $I_\zeta = \{p^m : p \text{ prime}, m \geq 1\}$ &
SHA-256 hash space: semigroup under data concatenation, no natural
inverse, provides unique digests (analogous to unique prime
factorisation). \\[0.4ex]

Hilbert-P\'olya group
  $I_{Z_\Gamma} = \Gamma \subset \PSL_2(\mathbb{R})$ &
ECDSA signature group: $E(\mathbb{F}_p)$-points on
$\mathrm{secp256k1}$ with chord-and-tangent addition, inverses
explicit, Casimir-like curve invariant $y^2 = x^3 + 7$. \\[0.4ex]

Weil quadratic form $Q_X$ (branch-transversal architecture; provides
gap positivity regardless of $\HBL$-status) &
Nakamoto consensus / Proof-of-Work: a protocol-level function that
validates the chain uniformly across hash-side and signature-side
participants, independently of which algebraic structure an
individual transaction uses. \\[0.4ex]

UCU: strict convexity of the governing functional yields a unique
minimiser (Section~\ref{sec:ucu}) &
Longest-chain rule: strictly increasing difficulty along the heaviest
chain yields a unique consensus head; cumulative proof-of-work is
the strictly convex ``score'' that selects it. \\[0.4ex]

Theorem~\ref{thm:NE-B}: no universal commuting operator for the
prime-shift family at tested Galerkin truncations &
Decentralisation: no single Bitcoin node holds a universal validating
key; the collective (primes / miners) validates by independent
contribution, not by a shared secret. \\[0.4ex]

Riemann hypothesis: all non-trivial zeros on the critical line &
Chain consistency: all accepted blocks lie on the canonical
(longest, Proof-of-Work-heaviest) chain. \\[0.4ex]

Three-component decomposition
  $\mathrm{RH}_X = \mathrm{GBC}_X + C^{(2)}_X + \mathrm{Hurwitz}_X$
  \cite{Geiger2026RHII} &
Three-layer protocol stack: transaction layer (GBC-analog, spectral
content of the operator), consensus layer ($C^{(2)}$-analog, Connes
approximation), historical persistence (Hurwitz-analog, reality of
zeros preserves the chain past re-orgs). \\[0.4ex]

Conjecture~\ref{conj:sge}: $\HBL$ is a binary algebraic invariant &
Protocol hybridism: Bitcoin's functioning requires the cooperation
of one semigroup-side component (hashes) and one group-side
component (signatures). Each is individually insufficient; their
branch-transversal coupling (the protocol itself) is the analog of
the Weil quadratic form. \\
\bottomrule
\end{tabular}
\end{center}

\subsection{Reading the trinity cryptographically}
\label{sec:blockchain-trinity}

The three meta-principles of Section~\ref{sec:trinity} receive a
compact cryptographic reading:

\begin{itemize}
\item \textbf{UCU as Proof-of-Uniqueness.}
The strict convexity of the governing functional is the protocol-level
analog of Bitcoin's longest-chain rule. In both cases, uniqueness of
the selected state is guaranteed by a monotone scoring function
(difficulty-accumulated work for Bitcoin; the Weil functional for
the zeta-family). Non-uniqueness would result in a fork; the
convexity condition prevents forks in both domains.

\item \textbf{SGE as protocol hybridism.}
Conjecture~\ref{conj:sge} asserts that each zeta-family lies
definitively on one side of the algebraic dichotomy. Bitcoin, by
contrast, is a single protocol that employs both sides: its hash
layer is $\HBL = \mathrm{NO}$-structured, its signature layer is
$\HBL = \mathrm{YES}$-structured. In this reading, Bitcoin is the
prototype of a \emph{hybrid protocol}---an object whose overall
architecture is impossible for a single zeta-family but natural when
the semigroup and group sides are coupled via a shared
branch-transversal mechanism.

\item \textbf{Weil-form as branch-transversal consensus.}
The Weil quadratic form $Q_X$ functions uniformly for
$\HBL = \mathrm{YES}$ and $\HBL = \mathrm{NO}$ families. Nakamoto
consensus similarly functions uniformly for hash-side operations
and signature-side operations. In both cases, the branch-transversal
layer is what allows the two algebraic worlds to cooperate
meaningfully.
\end{itemize}

\subsection{Why the correspondence is a controlled analogy}
\label{sec:blockchain-substance}

Three observations make the analogy structurally useful without
turning it into a theorem:

\begin{enumerate}
\item \emph{The dichotomy is precise.} Bitcoin's hash/signature
architecture combines a one-way digesting function (semigroup-like)
with an invertible authentication function (group-like). This mirrors,
but does not prove, the binary classification asserted by
Conjecture~\ref{conj:sge}.

\item \emph{Consensus is the algebra.} Nakamoto's insight was that
proof-of-work consensus can function as a branch-transversal layer
uniting the two sides. The Weil quadratic form plays the same role
in the zeta-family tree: it is the common arithmetic on which both
$\HBL$-regimes operate.

\item \emph{The uniqueness theorem is shared.} Bitcoin's longest-chain
rule is \emph{mathematically} a strictly convex selection principle;
the Weil functional's strict convexity is the same principle in the
continuous setting. Both are instances of a general convexity-implies-uniqueness
argument (UCU).
\end{enumerate}

The correspondence, accordingly, is best read as a structural echo.
The same trinity---UCU, SGE, Weil-form transversality---organises the
classification of zeta-type programmes and provides a compact
communication analogy for the architecture of the Bitcoin blockchain.

\subsection{Rigor disclaimer and further reading}
\label{sec:blockchain-disclaimer}

The present section is a \emph{speculative bridge}. It introduces
no new mathematical theorem and no new cryptographic result; the
correspondences above are conceptual pairings, not identities.
What the section does provide is a consolidated register in which
the trinity's structural content can be communicated to readers
from cryptography, computer science, and decentralised-systems
research. For a fuller informal development of the blockchain
analogy in the broader tree-of-zeta-families context (including
related motifs on primes as holographic boundary data and on time
as information accumulation), see the companion notes
\cite{MetaGalerkinBridge}, Appendix~I8--I10, and the supplementary
working notes referenced there.

% ============================================================
\section{Discussion and outlook}
\label{sec:outlook}

\subsection{Consequences for Dirichlet \texorpdfstring{$L$}{L}-functions}

Session-8 work on Dirichlet $L$-functions \cite{PaleyWienerDirichlet}
established that Paley-Wiener transfer via phenomenological kernels
(Gaussian or Hadamard-weighted pole) does not capture the
character-specific gap signatures. The boundary theorem
(Theorem~\ref{thm:boundary}) interprets this result: for low-conductor
real Dirichlet characters, no natural HP-BL operator is constructed
by simple perturbation of the Riemann case, and the observed
character-specific gaps must arise from the branch-transversal
Weil-form architecture, whose leading-order reductions have been
shown to fail. A dedicated Dirichlet-L atlas paper
\cite{DirichletAtlas} documents this failure in detail.

Under Conjecture~\ref{conj:sge}, this outcome is structurally
predicted: the Dirichlet transfer algebra is a semigroup, so
$\HBL(L(\cdot, \chi)) = \mathrm{NO}$ and no natural HP operator
can exist. The Dirichlet atlas is therefore not a failure of the
programme but a confirmation of its architectural boundaries.

A post-Zookeeper refinement is now part of the CORE status. The
Spectral Zookeeper \cite{Zookeeper2026} supplies an independent kernel
package for the Riemann cell: rank-one decomposition, Cauchy
interlacing, three-lemma closure, and spectral microclusters. This
does not change the boundary theorem; it refines the UBA layer. The
Riemann cell now carries two orthogonal kernel packages: v2.1 even
dominance as the published C2-sensitive layer, and the Zookeeper
three-lemma package as the theorem-level closure of the Riemann
$C^{(2)}$/MS2 step in the CCM/Fourier branch. For the Dirichlet cell
the natural back-transfer is Route~D:
a $\chi$-twisted CCM operator with a conductor rank-one term and a
$\chi$-weighted arithmetic perturbation. Thus the Dirichlet atlas
should be read as the first concrete UBA stress test of this transfer,
not merely as a negative Paley-Wiener numerical result.

\subsection{Proposed rigor hierarchy}
\label{sec:rigor}

Following the four-class hierarchy advocated in the tree meta-analysis
(theorem-level / conditional / diagnostic / speculative), each
component of the branch-locality programme can be classified as
follows:

\begin{center}
\begin{tabular}{@{}lll@{}}
\toprule
Component & Class & Source \\
\midrule
Definition~\ref{def:hp-bl} & theorem-level & this paper \\
Theorem~\ref{thm:NE-A} & theorem-level & \cite{Geiger2026RHII} \\
Theorem~\ref{thm:NE-B} ($N\leq 15$) & theorem-level & \cite{Geiger2026RHII} \\
Conjecture~\ref{conj:NE-B-full} & conditional & \cite{Geiger2026RHII} \\
Theorem~\ref{thm:selberg-hp} & theorem-level & \cite{Selberg1956} \\
Props.~\ref{prop:weyl-obstruction}--\ref{prop:fredholm-orientation} & diagnostic & \cite{Geiger2026RHIII}, \cite{MobiusGradedPrimeHub} \\
Theorem~\ref{thm:boundary} & conditional non-constancy theorem${}^*$ & synthesis \\
Conjecture~\ref{conj:sge} (SGE) & conjectural (finite-stress-tested) & this paper, \cite{SGENote} \\
\S\ref{sec:blockchain} (Blockchain reading) & speculative bridge & this paper \\
\bottomrule
\end{tabular}
\end{center}
${}^*$ conditional on Conjecture~\ref{conj:NE-B-full}.

\subsection{Transfer-data guardrail ledger}
\label{sec:transfer-data-ledger}

The preceding tables separate theorem-facing rows, calibration rows,
and speculative bridges. For reviewer navigation, the following ledger
makes the same separation operational: each data block is tied to the
claim it can support, the claim ceiling it cannot cross, and the section
where the detailed argument appears.

\begin{table}[htbp]
\centering
\scriptsize
\renewcommand{\arraystretch}{1.13}
\setlength{\tabcolsep}{3pt}
\begin{tabular}{@{}p{0.20\linewidth}p{0.24\linewidth}p{0.26\linewidth}p{0.16\linewidth}@{}}
\toprule
Data object & Supports & Claim ceiling & Location \\
\midrule
Canonical boundary rows:
Riemann, Selberg, Prime-Hub &
Non-constancy of branch-local Hilbert--P\'olya status &
Conditional non-constancy theorem, not an RH proof &
\S\ref{sec:boundary} \\
UCU instance table &
Reusable convexity template across analytic, physical, and algebraic cells &
Template identity only; no domain theorem without the local hypotheses &
\S\ref{sec:ucu} \\
Finite SGE stress tests:
Dedekind and Ihara-Petersen &
Calibration of the semigroup/group dichotomy &
Finite evidence or proxy; no $N\to\infty$ operator claim &
\S\ref{sec:sge-dedekind-test}, \S\ref{sec:sge-ihara-test} \\
Dirichlet and Zookeeper stress layer &
Branch-transversal tension and independent Riemann-cell kernel package &
Companion evidence only; no upgrade of Theorem~\ref{thm:boundary} &
\S\ref{sec:outlook} \\
Blockchain reading &
Controlled analogy for ledgers, proof-of-work, and boundary witnesses &
Speculative bridge, not mathematical evidence &
\S\ref{sec:blockchain} \\
Boundary-taxonomy review fields &
Future audit fields for zoo population, two-collar cases, and finite-pattern checks &
Review checklist; finite residue patterns and two-collar passes remain
diagnostic until a Sturmian/irrational-rotation, homotopy-winding, or
genuine limit argument is supplied &
\S\ref{sec:outlook} \\
\bottomrule
\end{tabular}
\caption{Transfer-data guardrail ledger. The table records what each
data block is allowed to support and where it must stop.}
\label{tab:transfer-data-ledger}
\end{table}

\subsection{Future directions}

\begin{enumerate}
\item \textbf{Formal proof of Conjecture~\ref{conj:NE-B-full}.} An
  SVD-density argument for $N = \infty$ is the single most valuable
  next step: it would upgrade Theorem~\ref{thm:boundary} from
  conditional to unconditional.
\item \textbf{Extend the Dedekind-zeta evidence.} The finite
  $\mathbb{Q}(\sqrt{-5})$ stress test supports the SGE prediction but
  does not replace an $N\to\infty$ SVD-density argument. The next step is
  to extend the test beyond this single finite certificate, across
  additional number fields and toward a genuine infinite-dimensional
  limit statement.
\item \textbf{Ihara zeta and Yang-Lee roots.} Ihara's case is the
  graph-theoretic analogue of Selberg's and is expected by
  Conjecture~\ref{conj:sge} to satisfy $\HBL = \mathrm{YES}$, with
  the Bass-Hashimoto matrix serving as $H_X$. Yang-Lee zeros, by
  contrast, are partition-function roots parameterised by
  temperature and should on present evidence sit in the flow-zeta
  class of Section~\ref{sec:three-class-taxonomy}, as a
  statistical-mechanical relative of CRM.

\item \textbf{Population of the zoo.} The six zoo entries exhibited
  so far (Riemann, Selberg, Prime-Hub, CRM, Dedekind
  $\mathbb{Q}(\sqrt{-5})$, Ihara-Petersen) plus the Dirichlet atlas
  companion \cite{DirichletAtlas} are illustrative, not exhaustive.
  A scan of the surrounding literature and the companion tree
  documents \cite{CandidatesMasterList} identifies eight further
  families that fit the three-class taxonomy cleanly and sharpen
  the SGE dichotomy:
  \begin{itemize}
  \item \textbf{Automorphic $L(\pi, s)$ for $\mathrm{GL}_n$}
  (Langlands programme): geometric/representation-theoretic,
  \textbf{candidate SGE-YES}. The transfer algebra is the Hecke
  algebra, and the ambient reductive group supplies Casimir and
  trace-formula provenance. This makes automorphic $L$-functions a
  natural YES-side candidate class \cite{Langlands1970,CogdellPS1994},
  not yet a settled
  Hilbert--P\'olya operator for their zero ordinates in the sense of
  Definition~\ref{def:hp-bl}.
  \item \textbf{Hecke $L(s, \lambda)$ over number fields}:
  arithmetic, \textbf{SGE-NO}. Sister species to Dirichlet and
  Dedekind---prime-ideal semigroup plus character twist
  \cite{IwaniecKowalski}.
  \item \textbf{Artin $L$-functions}: geometric-compact,
  \textbf{SGE-YES}. Transfer algebra is a finite Galois group;
  prototype of the Galois-driven YES-side
  \cite{Artin1931,Chebotarev1926}.
  \item \textbf{Ruelle-Anosov zeta $\zeta_\varphi(s)$} for hyperbolic
  flows: flow-noncompact, \textbf{SGE-YES}. Direct generalisation
  of CRM beyond the $\PSL_2(\mathbb{R})$-boost case; see
  \cite{DyatlovZworski2016,DyatlovZworski2019}.
  \item \textbf{Selberg on the non-compact quotient}
  $\mathrm{PSL}_2(\mathbb{R})/\mathrm{SO}(2) = \mathbb{H}$:
  flow-noncompact, \textbf{SGE-YES}. Continuous spectrum plus
  Eisenstein series. Algebraically almost identical to CRM,
  dynamically distinct.
  \item \textbf{Black-hole quasinormal-mode spectra}
  (\cite{DeClerckHartnollYang2025}): flow-noncompact,
  \textbf{SGE-YES}. Scattering resonances on asymptotically
  de~Sitter backgrounds; the Wheeler--DeWitt wavefunctions of 5d
  BKL dynamics exhibit automorphic-$L$ and complex-primon-gas
  structures linking zeta-like spectra to quantum-gravity regimes.
  The first flow-zeta with realistic hopes of experimental access
  via gravitational-wave ringdown observations (LIGO/Virgo/Einstein
  Telescope).
  \item \textbf{Yang-Lee zeros}:\\[-0.6ex]
  ferromagnetic partition functions;
  conjecturally flow-noncompact, \textbf{SGE-YES}. Temperature as
  continuous parameter; the Lee-Yang circle theorem plays the role
  of an RH-analog. Statistical-mechanical cousin of CRM.
  \item \textbf{Bost-Connes system} (1995): a quantum-statistical
  realisation of $\zeta(s)$ as partition function with spontaneous
  symmetry breaking at $\beta = 1$. Sits on the bridge between the
  arithmetic and flow classes and connects the zoo directly to
  Connes' adelic framework \cite{BostConnes1995,Connes1999, Connes2026}.
  \end{itemize}
  Each of these deserves a dedicated FST-Mathematics supplement in
  due course. A second tier of twelve further constructions
  (Epstein, Witten-zeta, multiple-zeta-values, Hurwitz, Lerch,
  Arakelov, motivic $L$, chiral-zeta, RG-zetas, Gromov-Witten,
  Reidemeister torsion, QFT spectral zetas) and a separate column
  of biological Pattern-A-derived zetas (Kauffman-attractor,
  protein-folding spectral, replicator Jacobian) are catalogued in
  the companion tree document \cite{CandidatesMasterList}.

\item \textbf{Biological Pattern-A instances as bridges.} The
  FST-Biology supplement programme \cite{FSTBiology} develops
  Pattern-A in molecular, chemical, and behavioural domains
  (protein folding as a Nash equilibrium, gene regulatory network
  attractors, immune-system coevolution). These are primarily
  residents of the physical side of the programme, but generate
  zeta-like counting functions of their own (Kauffman-attractor
  zetas for Boolean gene networks, protein-folding spectral zetas,
  replicator Jacobian spectra) that could in principle populate
  dedicated enclosures in the zoo. Their primary classification
  is on the physical side of FST; their mathematical structure is
  listed in the companion candidate document.

\item \textbf{Adelic upgrade: Riemann revisited.} Connes' adelic
  framework \cite{Connes1999, Connes2026} replaces the classical
  prime-power semigroup $\{p^m\}$ with the idele class group
  $\mathbb{A}_K/K^\times$, which \emph{is} a group. If that
  programme closes, Riemann migrates from the SGE-NO column
  (classical prime-shift transfer algebra) to the SGE-YES column
  (adelic transfer algebra), without contradicting
  Theorem~\ref{thm:boundary}---the SGE conjecture describes
  precisely what a successful Hilbert-P\'olya programme must
  achieve.

\item \textbf{Explicit Prime-Hub construction.} Any future
  construction must explicitly navigate
  Propositions~\ref{prop:weyl-obstruction}--\ref{prop:fredholm-orientation}.
  Under Conjecture~\ref{conj:sge}, the construction of
  $G_\mathrm{prime}$ succeeds in producing a Hilbert-P\'olya
  operator if and only if the resulting cycle-concatenation
  algebra is a group. In addition, OP5.3 and OP5.5 must be separated
  into an Euler/M\"obius channel and a global spectral channel, or
  encoded by a genuine superdeterminant; a naked positive Fredholm
  determinant has the wrong orientation at zeta zeros.

\item \textbf{Boundary-taxonomy ledger.} Recent transfer checks turn
  the zoo from a list of species into a guardrail ledger. Four boundary
  types, now extended by the Waisen/Ihara two-collar and Sturmian
  finite-pattern gate, should be recorded for future entries:
  \begin{center}
  \scriptsize
  \begin{tabular}{@{}p{0.20\textwidth}p{0.32\textwidth}p{0.32\textwidth}@{}}
  \toprule
  Type & Test & Guardrail \\
  \midrule
  Prunable appendage &
  Core-prune Ihara-like side corridors before classification &
  Finite tree legs that leave the zeta/operator object invariant are not new species \\
  Schur-cofactor obstruction &
  For $G=(\begin{smallmatrix}A&b\\ b^t&c\end{smallmatrix})$, require full-rank $A$ and check $s=c-b^tA^{-1}b$, $Q_B=\beta^2/s$ &
  $\det(G)\neq0$ is only a shortcut with full-rank source block and non-zero coupling \\
  Fixed target vs global transfer &
  Separate fixed target probes from growing windows and limit statements &
  Hitting selected targets is not a global Hilbert--P\'olya or SGE transfer \\
  Local pass vs global section &
  Assemble local root, branch or verifier passes into a coherent low global section &
  Local passes without a section remain diagnostic \\
  Core collar vs boundary collar &
  Separate Schur-supported core movement from boundary-circle claims; record support set and support rank &
  Local defects move only the rank-visible cluster; boundary claims need a Rouch\'e/argument-principle window with homotopy and winding checks \\
  Finite periodic pattern &
  Audit residue-class or period laws against phase, Sturmian/irrational-rotation, or true limit arguments &
  Long finite stability is diagnostic only; the Waisen S3 warning case breaks the naive pattern at $L=46$ \\
  \bottomrule
  \end{tabular}
  \end{center}
  Suggested review fields include
  \path|survives_core_pruning|, \path|schur_scalar|,
  \path|fixed_target_probe|, \path|growth_window_status|,
  \path|global_limit_status|, \path|selection_rule|,
  \path|collar_type|, \path|support_set|,
  \path|support_rank|, \path|fixed_multiplicity_lower_bound|,
  \path|rouche_window_status|,
  \path|homotopy_winding_checked|, and
  \path|finite_pattern_status|.
  This ledger does not change Theorem~\ref{thm:boundary}; it makes
  future zoo population harder to overread from local or finite
  certificates.

\item \textbf{Independent replication and a reviewer-verifiability
  index.} The trilogy \cite{Geiger2026RHII,Geiger2026RHIII}
  currently comprises \textasciitilde{}82 PDF pages, 76 theorem
  objects, an eight-step reduction chain, 33 CAP certificates, and
  twenty Python scripts. The critique review
  \cite{RHCritiqueReview2026}, while finding no fatal gap,
  identifies documentation-ergonomic weaknesses (K6): absence of a
  central symbol glossary, no dependency graph A1--A8, no unified
  \texttt{reproduce.sh} pipeline, no ``reviewer fast path'' note.
  We propose a \emph{Reviewer Kit} companion Zenodo record
  containing (a) a symbol glossary, (b) a dependency DAG linking
  lemmas to scripts to A-steps, (c) a
  \texttt{scripts/REPRODUCE.md} with one-liner invocations and
  expected runtimes, (d) a Lean4 skeleton for the Shift Parity
  Lemma---its purely algebraic structure ($N \geq 3$, closed-form
  determinant and trace) makes it the natural first target for
  serious formalisation, and (e) a two-page ``reviewer fast path''
  stating how each A-step can be independently verified in a
  bounded time budget. An accompanying \emph{reviewer-verifiability
  index} per A-step (analytic / CAP / hybrid, estimated person-hours
  for independent verification) would make the trilogy's
  replication cost explicit. Crucially, this programme requires no
  change to the trilogy itself; it is an additive companion that
  addresses K6 without touching the proof.
\end{enumerate}

% ============================================================
\section*{AI Disclosure}

This manuscript was prepared with extensive AI-assisted support for
literature organisation, drafting, translation, consistency checks,
LaTeX maintenance, and mathematical review prompts. AI systems were used
as tools and are not authors, co-authors, acknowledgement recipients,
truth authorities, or independent validators. All mathematical claims,
interpretations, and publication decisions remain the sole responsibility
of the author.

% ============================================================
\begin{thebibliography}{99}

% --- Internal project references ---

\bibitem{Geiger2026RHI}
L.~Geiger, \emph{The Riemann Hypothesis via the v2.0 Method,
Part I: Foundations and Obstructions}, Preprint (2026), Zenodo concept DOI
\href{https://doi.org/10.5281/zenodo.19035640}{10.5281/zenodo.19035640}.

\bibitem{MetaGalerkinBridge}
L.~Geiger, \emph{A Meta Galerkin Bridge: Diagnostic Decomposition
of RH-Type Programmes}, Working draft (2026),
\texttt{paper/META\_GALERKIN\_BRIDGE.tex} in the archived predecessor
tree. Appendix I collects
the ten essay-level motivations referenced here, including the
blockchain-themed items~I8--I10.

\bibitem{RFEPFramework}
L.~Geiger, \emph{Functional Stability Theory --- Physical
Instantiation via the Renormalized Free Energy Principle},
Preprint, companion paper (2026), Zenodo concept DOI
\href{https://doi.org/10.5281/zenodo.19036190}{10.5281/zenodo.19036190};
public repository
\href{https://github.com/research-line/functional-stability-theory/tree/main/masters/spectrum-duality}{research-line/functional-stability-theory/masters/spectrum-duality}.
The current archival release is v1.9 under the RFEP/Spectrum-Duality title
(Record DOI
\href{https://doi.org/10.5281/zenodo.20758955}{10.5281/zenodo.20758955};
Concept DOI
\href{https://doi.org/10.5281/zenodo.19036190}{10.5281/zenodo.19036190}).
The public repository path carries a local v1.10 candidate; v1.10 is not
published. ``RFEP v1.8'' is retained only as the internal
terminology/update-plan label in the source tree, not as a separate archival
release.

\bibitem{FSTYangMills}
L.~Geiger, \emph{FST-Physics: Yang--Mills Mass Gap via Pattern A},
Preprint (2025/2026), Zenodo DOI
\href{https://doi.org/10.5281/zenodo.19087433}{10.5281/zenodo.19087433}.

\bibitem{FSTTurbulence}
L.~Geiger, \emph{FST-Physics: Kolmogorov K41 Spectrum via
Pattern A}, Preprint (2025/2026), Zenodo DOI
\href{https://doi.org/10.5281/zenodo.19056813}{10.5281/zenodo.19056813}.

\bibitem{Geiger2026RHII}
L.~Geiger, \emph{A Conditional Reduction of the Riemann Hypothesis to
Even Dominance of the Weil Quadratic Form}, Preprint (2026),
Zenodo concept DOI
\href{https://doi.org/10.5281/zenodo.19035640}{10.5281/zenodo.19035640}.

\bibitem{Geiger2026RHIII}
L.~Geiger, \emph{The Riemann Hypothesis via the v2.0 Method,
Part III: Conclusio and Open Problems}, Preprint (2026),
Zenodo concept DOI
\href{https://doi.org/10.5281/zenodo.19035640}{10.5281/zenodo.19035640}.

\bibitem{DirichletAtlas}
L.~Geiger, \emph{A Weil-Kernel Numerical Atlas for Dirichlet
$L$-Function Sector Gaps: Galerkin Diagnostics and the Boundaries of
Leading-Order Theory}, Preprint (2026), Zenodo concept DOI
\href{https://doi.org/10.5281/zenodo.19960809}{10.5281/zenodo.19960809};
source path \texttt{CORE/zoo-mapping/paper/DIRICHLET\_CHARACTER\_ATLAS\_v1\_en.tex}.

\bibitem{PaleyWienerDirichlet}
L.~Geiger, \emph{Paley-Wiener Transfer for Dirichlet $L$-functions:
First-Order Perturbation and Numerical Validation}, Working notes
(2026),
\texttt{CORE/zoo-mapping/PALEY\_WIENER\_DIRICHLET\_TRANSFER.md}.

\bibitem{SGENote}
L.~Geiger, \emph{The Semigroup-Group Equivalence: A Unifying
Conjecture for Branch-Locality}, Working notes (2026),
\texttt{paper\_NE\_B\_BOUNDARY/SEMIGROUP\_GROUP\_EQUIVALENCE.md}.

\bibitem{DedekindNEBTest}
L.~Geiger, \emph{Dedekind NE-B Analog Test for
$\mathbb{Q}(\sqrt{-5})$: First Empirical Probe of the
Semigroup-Group Equivalence}, Computational notes (2026),
\texttt{paper\_NE\_B\_BOUNDARY/\_results/DEDEKIND\_NE\_B\_TEST.md}.

\bibitem{IharaPetersenTest}
L.~Geiger, \emph{Ihara-Zeta SGE YES-side Test for the Petersen
Graph: Bass-Hashimoto, Ramanujan, and the Laplacian as HP-Operator},
Computational notes (2026),
\texttt{paper\_NE\_B\_BOUNDARY/\_results/IHARA\_PETERSEN\_SGE.md}.

\bibitem{MobiusGradedPrimeHub}
L.~Geiger, \emph{M\"obius-Graded Prime-Hub: Determinant Orientation,
Exterior-Fock Normal Form, and the OP5 Two-Channel Split},
Computational note (2026),
\texttt{CORE/zetazoo/\_results/MOBIUS\_GRADED\_PRIMEHUB\_TOY.md}.

\bibitem{PrimeWeilBridgeDefect}
L.~Geiger, \emph{Prime--Weil Bridge Defect: Diagnostic and
Null-Free Fourier-CCM Operator Controls}, Computational artefact
(2026),
\texttt{CORE/zetazoo/\_results/PRIME\_WEIL\_BRIDGE\_DEFECT.md} and
\texttt{CORE/zetazoo/\_results/PRIME\_WEIL\_BRIDGE\_DEFECT\_OPERATOR\_FOURIER.md}.

\bibitem{CRMFlowZetaNote}
L.~Geiger, \emph{CRM as Flow-Zeta: Reformulation of the Cooperative
Renormalization Model of Dark Energy as a Ruelle-Dyatlov-Zworski
Zeta on the Hyperbolic Plane}, Working note (2026),
\texttt{04\_FST\_COSMOLOGY/CRM\_AS\_FLOW\_ZETA.md}.

\bibitem{Zookeeper2026}
L.~Geiger, \emph{The Spectral Zookeeper: The Riemann Hypothesis via
Spectral Microcluster Closure in the Connes--Consani--Moscovici
Framework}, Preprint (2026), Zenodo DOI
\href{https://doi.org/10.5281/zenodo.19673126}{10.5281/zenodo.19673126}.

\bibitem{CRMV}
L.~Geiger, \emph{Cooperative Renormalization Model of Dark Energy V:
The Saturation Theorem}, Preprint (2026), Zenodo concept DOI
\href{https://doi.org/10.5281/zenodo.18728935}{10.5281/zenodo.18728935}.

\bibitem{DyatlovZworski2019}
S.~Dyatlov and M.~Zworski, \emph{Mathematical Theory of Scattering
Resonances}, Graduate Studies in Mathematics \textbf{200}, American
Mathematical Society, Providence, RI (2019),
DOI~\href{https://doi.org/10.1090/gsm/200}{10.1090/gsm/200}.

\bibitem{CandidatesMasterList}
L.~Geiger, \emph{Zeta Zoo Candidates --- Master List}, Companion
tree document (2026),
\texttt{CORE/CANDIDATES\_MASTER\_LIST.md} in the
\texttt{.ZETA-ZOO} project tree. Maintains the full
catalogue of zeta-type family candidates identified across RH
v2.1, EFP-Taxonomy, V2-Physics-Testcases, FST-Biology, and
supplementary Ideenanhang sources.

\bibitem{FSTBiology}
L.~Geiger, \emph{Functional Stability Theory --- Biological
Pattern-A Instances (FST-I Thermodynamic, FST-II Chemical, FST-III
Biological Stability)}, Preprint series (2026),
Zenodo concept DOIs
\href{https://doi.org/10.5281/zenodo.20130544}{10.5281/zenodo.20130544},
\href{https://doi.org/10.5281/zenodo.20130563}{10.5281/zenodo.20130563},
and \href{https://doi.org/10.5281/zenodo.20130573}{10.5281/zenodo.20130573}.
Three-paper trilogy on scale-invariant Pattern-A realisations from
molecular to behavioural biology.

% --- External classical references (verified via WebSearch 2026-04-16) ---

\bibitem{Polya1982letter}
G.~P\'olya, letter to A.~M.~Odlyzko, 3~January 1982, recounting a
Göttingen conversation with E.~Landau (ca.\ 1912--1914) about a
possible operator-theoretic explanation for the Riemann hypothesis.
Preserved by Odlyzko at
\url{https://www-users.cse.umn.edu/~odlyzko/polya/}.

\bibitem{Montgomery1973}
H.~L.~Montgomery, \emph{The pair correlation of zeros of the zeta
function}, in: \emph{Analytic Number Theory} (Proc.\ Sympos.\ Pure
Math.\ \textbf{XXIV}, St.\ Louis 1972), pp.~181--193, American
Mathematical Society, Providence, RI (1973). Earliest published
mention of the Hilbert--P\'olya suggestion as folklore.

\bibitem{Berry1985}
M.~V.~Berry, \emph{Semiclassical theory of spectral rigidity},
Proc.\ Roy.\ Soc.\ London A \textbf{400}, 229--251 (1985), DOI
\href{https://doi.org/10.1098/rspa.1985.0078}{10.1098/rspa.1985.0078}.

\bibitem{BerryKeating1999}
M.~V.~Berry and J.~P.~Keating, \emph{The Riemann zeros and
eigenvalue asymptotics}, SIAM Rev.\ \textbf{41} (2), 236--266 (1999),
DOI
\href{https://doi.org/10.1137/S0036144598347497}{10.1137/S0036144598347497}.

\bibitem{Connes1999}
A.~Connes, \emph{Trace formula in noncommutative geometry and the
zeros of the Riemann zeta function}, Selecta Math.\ (New Series)
\textbf{5} (1), 29--106 (1999), DOI
\href{https://doi.org/10.1007/s000290050042}{10.1007/s000290050042}.

\bibitem{Connes2026}
A.~Connes, \emph{The Riemann Hypothesis: Past, Present and a
Letter Through Time}, arXiv preprint (2026), arXiv:2602.04022
[math.NT], \url{https://arxiv.org/abs/2602.04022}.

\bibitem{ConnesConsaniMoscovici2025}
A.~Connes, C.~Consani, and H.~Moscovici, \emph{Zeta Spectral
Triples}, arXiv preprint (2025), arXiv:2511.22755 [math.NT],
\url{https://arxiv.org/abs/2511.22755}. Companion paper to
\cite{Connes2026}.

\bibitem{GeigerObstructionClassifiers}
L.~Geiger, \emph{Obstruction Classifiers for RH-Type Problems},
Working paper (2026), cf.\ \texttt{\_archive/old\_obstruction\_classifier\_wrapper/Obstruction\_Classifiers\_v1\_en.tex}
in the archived predecessor tree.

\bibitem{SGEControlExperiment2026}
L.~Geiger, \emph{SGE Control Experiment --- Discriminating Test},
Working note and computational artefact (2026),
\texttt{CORE/zetazoo/\_results/SGE\_CONTROL\_EXPERIMENT.md},
produced by
\texttt{CORE/zetazoo/\_scripts/sge\_control\_experiment.py}.
Server-certified run (ellmos-services, 2026-04-17) reporting
centraliser-dimension scaling for cyclic $\mathbb{Z}/N$,
elementary abelian $(\mathbb{Z}/2)^k$ and random-symmetric
baselines.

\bibitem{RHCritiqueReview2026}
L.~Geiger, \emph{RH v2.1 --- Systematic Critique Review (K1--K6)},
Working note (2026), \texttt{SESSIONS/\_reviews/RH\_CRITIQUE\_REVIEW\_2026-04-17.md}
in the \texttt{.ZETA-ZOO} project tree.
Four-agent parallel assessment of the Foundation--Even-Dominance--Conclusio
trilogy against six external critiques; verdict: all partially valid,
none fatal.

\bibitem{EkelandTemam1999}
I.~Ekeland and R.~T\'emam, \emph{Convex Analysis and Variational
Problems}, SIAM Classics in Applied Mathematics \textbf{28}, Society
for Industrial and Applied Mathematics, Philadelphia (1999),
ISBN 978-0-89871-450-0. Standard reference for Mazur's lemma and
the strict-convexity-plus-coercivity uniqueness principle used in
Lemma~\ref{lem:ucu}.

% --- Added external classical / contemporary references (WebSearch-verified 2026-04-17) ---

\bibitem{KeatingSnaith2000}
J.~P.~Keating and N.~C.~Snaith, \emph{Random Matrix Theory and
$\zeta(1/2 + it)$}, Commun.\ Math.\ Phys.\ \textbf{214} (2000),
pp.~57--89, DOI~\href{https://doi.org/10.1007/s002200000261}{10.1007/s002200000261}.

\bibitem{Bombieri2000}
E.~Bombieri, \emph{Problems of the Millennium: The Riemann
Hypothesis}, Clay Mathematics Institute Official Problem
Description (2000),
\url{https://www.claymath.org/wp-content/uploads/2022/05/riemann.pdf}.

\bibitem{Conrey2003}
J.~B.~Conrey, \emph{The Riemann Hypothesis}, Notices of the AMS
\textbf{50} (2003), no.~3, pp.~341--353,
\url{https://www.ams.org/notices/200303/fea-conrey-web.pdf}.

\bibitem{BostConnes1995}
J.-B.~Bost and A.~Connes, \emph{Hecke algebras, type III factors
and phase transitions with spontaneous symmetry breaking in number
theory}, Selecta Math.\ (New Series) \textbf{1} (1995), no.~3,
pp.~411--457,
DOI~\href{https://doi.org/10.1007/BF01589495}{10.1007/BF01589495}.

\bibitem{Langlands1970}
R.~P.~Langlands, \emph{Problems in the theory of automorphic
forms}, in: C.~T.~Taam (ed.), \emph{Lectures in Modern Analysis
and Applications III}, Lecture Notes in Mathematics \textbf{170},
Springer, Berlin--Heidelberg (1970), pp.~18--61,
DOI~\href{https://doi.org/10.1007/BFb0079065}{10.1007/BFb0079065}.

\bibitem{DeClerckHartnollYang2025}
M.~De~Clerck, S.~A.~Hartnoll, and M.~Yang,
\emph{Wheeler--DeWitt wavefunctions for 5d BKL dynamics, automorphic
$L$-functions and complex primon gases}, arXiv preprint (2025),
arXiv:2507.08788 [hep-th],
\url{https://arxiv.org/abs/2507.08788}.

\bibitem{Artin1931}
E.~Artin, \emph{Zur Theorie der L-Reihen mit allgemeinen
Gruppencharakteren}, Abh.\ Math.\ Sem.\ Univ.\ Hamburg \textbf{8}
(1931), pp.~292--306,
DOI~\href{https://doi.org/10.1007/BF02941010}{10.1007/BF02941010}.

\bibitem{Chebotarev1926}
N.~Tschebotareff, \emph{Die Bestimmung der Dichtigkeit einer Menge
von Primzahlen, welche zu einer gegebenen Substitutionsklasse
geh\"oren}, Math.\ Ann.\ \textbf{95} (1926), pp.~191--228,
DOI~\href{https://doi.org/10.1007/BF01206606}{10.1007/BF01206606}.

\bibitem{CogdellPS1994}
J.~W.~Cogdell and I.~I.~Piatetski-Shapiro, \emph{Converse theorems
for $GL_n$}, Publ.\ Math.\ Inst.\ Hautes \'Etudes Sci.\ \textbf{79}
(1994), pp.~157--214,
\url{https://www.numdam.org/item/PMIHES_1994__79__157_0/}.

\bibitem{DyatlovZworski2016}
S.~Dyatlov and M.~Zworski, \emph{Dynamical zeta functions for
Anosov flows via microlocal analysis}, Ann.\ Sci.\ \'Ec.\ Norm.\
Sup\'er.\ (4) \textbf{49} (2016), no.~3, pp.~543--577,
DOI~\href{https://doi.org/10.24033/asens.2290}{10.24033/asens.2290},
arXiv:1306.4203.

\bibitem{Selberg1956}
A.~Selberg, \emph{Harmonic analysis and discontinuous groups in
weakly symmetric Riemannian spaces with applications to Dirichlet
series}, J.\ Indian Math.\ Soc.\ \textbf{20}, 47--87 (1956).

\bibitem{Marklof2008}
J.~Marklof, \emph{Selberg's trace formula: an introduction}, in
\emph{Hyperbolic Geometry and Applications in Quantum Chaos and
Cosmology}, Cambridge University Press (2011), pp.~83--119,
\href{https://doi.org/10.48550/arXiv.math/0407288}{doi:10.48550/arXiv.math/0407288}.

\bibitem{Weil1952}
A.~Weil, \emph{Sur les ``formules explicites'' de la th\'eorie des
nombres premiers}, Meddelanden fr\aa n Lunds Universitets
Matematiska Seminarium, Supplementband tillägnet Marcel Riesz,
pp.~252--265, C.~W.~K.\ Gleerup, Lund (1952).

\bibitem{IwaniecKowalski}
H.~Iwaniec and E.~Kowalski, \emph{Analytic Number Theory}, American
Mathematical Society Colloquium Publications \textbf{53}, AMS,
Providence, RI (2004), xi+615~pp., ISBN 0-8218-3633-1.

\bibitem{Titchmarsh1986}
E.~C.~Titchmarsh, \emph{The Theory of the Riemann Zeta-Function},
2nd edition, revised by D.~R.~Heath-Brown, Oxford University Press,
Oxford (1986), ISBN 0-19-853369-1.

\end{thebibliography}

\end{document}
